% !TeX program = pdflatex
% Zbior zadan do skryptu topologii algebraicznej, iteracja XIX --- audyt globalny.
% Zadania służą także odtworzeniu trudniejszych dowodów i rachunków.
\documentclass[11pt,a4paper,oneside,openany]{book}
\usepackage[T1]{fontenc}
\usepackage[utf8]{inputenc}
\usepackage[polish]{babel}
\usepackage{lmodern}
\usepackage{amsmath,amssymb,amsthm,mathtools}
\usepackage{microtype}
\usepackage{geometry}
\usepackage{enumitem}
\usepackage[unicode,colorlinks=true,linkcolor=blue!40!black]{hyperref}
\usepackage{xcolor}
\geometry{margin=2.55cm}
\pagestyle{plain}
\setlength{\parindent}{0pt}
\setlength{\emergencystretch}{2em}
\setlist[enumerate]{leftmargin=*,itemsep=.65em,topsep=.7em}
\newcommand{\F}{\mathbb F}
\newcommand{\R}{\mathbb R}
\newcommand{\C}{\mathbb C}
\newcommand{\Z}{\mathbb Z}
\newcommand{\Q}{\mathbb Q}
\newcommand{\N}{\mathbb N}
\newcommand{\T}{\mathbb T}
\newcommand{\dd}{\mathrm d}
\newcommand{\D}{\mathrm D}
\newcommand{\Hom}{\operatorname{Hom}}
\newcommand{\rank}{\operatorname{rank}}
\newcommand{\im}{\operatorname{im}}
\newcommand{\kerop}{\operatorname{ker}}
\newcommand{\Cl}{\operatorname{Cl}}
\newcommand{\Spin}{\operatorname{Spin}}
\newcommand{\Gr}{\operatorname{Gr}}
\newcommand{\Ric}{\operatorname{Ric}}
\newcommand{\Scal}{\operatorname{Scal}}
\newcommand{\Wh}{\operatorname{Wh}}
\newcommand{\id}{\operatorname{id}}
\newcommand{\wskazowka}[1]{%
  \par\smallskip\noindent{\footnotesize\color{black!75}\emph{Wskazówka.}\enspace #1}\par
}
\begin{document}
\frontmatter
\begin{center}
{\LARGE\bfseries Topologia algebraiczna: zadania i wskazówki}\par\vspace{.6em}
{\large Oddzielny zeszyt ćwiczeń do skryptu}\par
\end{center}
\bigskip
Zachowano zadania i wskazówki z wcześniejszego zestawu ćwiczeń.
Wprowadzono także zadania o fibracjach, Eilenbergu--Zilberze,
teorii \v{C}echa, dualności, podstawy algebry abstrakcyjnej, ciągów spektralnych, wież Postnikowa, operacji kohomologicznych i homotopii stabilnej, a także Pontriagina--Thoma, kobordyzmów, homologii Morse'a i twierdzenia Lefschetza.
\tableofcontents
\mainmatter
\chapter{Podstawowe konstrukcje topologiczne}
\begin{enumerate}
\item Udowodnij, że iloraz $[0,1]/(0\sim1)$ jest homeomorficzny z $S^1$.
\wskazowka{Użyj odwzorowania $t\mapsto e^{2\pi it}$ i zwartości.}
\item Wyznacz deformacyjny retrakt walca $S^1\times[0,1]$ na jeden z okręgów brzegowych. Czy istnieje analogiczna retrakcja dysku $D^2$ na brzeg?
\wskazowka{Dla drugiego pytania skorzystaj z grupy podstawowej.}
\item Wykaż, że zawieszenie $S^0$ jest homeomorficzne z $S^1$ i opisz obraz każdego odcinka $\{x\}\times[0,1]$.

\item \textbf{Stożek a retrakcja.} Zdefiniuj stożek $CX=(X\times[0,1])/(X\times\{1\})$ nad niepustą przestrzenią $X$ i skonstruuj homotopię ściągającą go do wierzchołka. Dlaczego nie wynika stąd, że $X$ jest ściągalna?
\wskazowka{Zachowaj współrzędną punktu $x$ i liniowo przesuwaj parametr ku $1$. Przestrzeń $X$ nie musi być retraktem stożka.}
\item \textbf{Iloraz a ciągłość.} Niech $q:X\to X/\sim$ będzie odwzorowaniem ilorazowym. Udowodnij, że $F:X/\sim\to Y$ jest ciągłe wtedy i tylko wtedy, gdy $Fq$ jest ciągłe. Zastosuj wynik do sklejenia dwóch końców odcinka.
\item \textbf{Zawieszenie sumy rozłącznej.} Porównaj $\Sigma(S^0\sqcup S^0)$ i $\Sigma(S^0\vee S^0)$; przed rachunkiem jawnie określ, co oznacza punkt bazowy w sumie klinowej.
\wskazowka{Nie myl sumy rozłącznej z sumą klinową; policz składowe i pętle w prostych modelach CW.}

\end{enumerate}
\chapter{Podstawy algebry abstrakcyjnej}
\begin{enumerate}
\item Niech $G=\Z/12\Z$ i $H=4\Z/12\Z$. Wyznacz $G/H$ oraz sprawdź pierwszy izomorfizm dla mapy $\Z/12\Z\to\Z/4\Z$.
\item Zbadaj działanie $S_4$ na zbiorze par nieuporządkowanych elementów $\{1,2,3,4\}$. Wyznacz orbitę i stabilizator $\{1,2\}$.
\wskazowka{Stabilizator może zamieniać elementy wewnątrz obu bloków.}
\item W grupie $\Z\rtimes_{(-1)}\Z$ oblicz $(3,2)(-1,3)$ i odwrotność $(3,2)$. Wykaż, że podgrupa $\Z\times\{0\}$ jest normalna i że cała grupa nie jest abelowa.
\item Udowodnij, że $\langle a,b\mid a^2=b^3=1\rangle$ ma homomorfizm na $S_3$ wysyłający $a$ na transpozycję, $b$ na cykl długości trzy. Czy przesądza to o izomorfizmie?
\wskazowka{Rozróżnij homomorfizm surjektywny od różnowartościowego.}
\item Uzasadnij, że $R=\Z[x]/(2,x)$ jest izomorficzny z $\Z/2\Z$. Zbadaj, czy ideał $(2,x)$ w $\Z[x]$ jest główny.
\item Dla $f:\Z^3\to\Z^2$, $f(a,b,c)=(a+2b,b+3c)$, oblicz $\ker f$ i $\operatorname{coker} f$. Czy wynik odpowiada pierwszemu twierdzeniu o izomorfizmie?
\item Rozstrzygnij, czy ciąg $0\to\Z\xrightarrow{\cdot m}\Z\to\Z/m\Z\to0$ rozszczepia się dla $m\ge2$.
\item Skonstruuj odwzorowanie $\Z/6\Z\times\Z/4\Z\to\Z/2\Z$, którego uniwersalność pokazuje, że $\Z/6\otimes\Z/4\cong\Z/2$. Następnie oblicz $\Z/6\otimes\Z/5$.
\item Wykaż wprost, że $\Z/2$ nie jest płaskim $\Z$-modułem i że $\Z$ jest projektywna.
\item Niech $A=\Z/8\Z$, $B=\Z/6\Z$. Wyznacz $A\otimes B$, $\operatorname{Tor}_1^{\Z}(A,B)$ oraz $\operatorname{Ext}^1_\Z(A,B)$.
\wskazowka{Dla $\operatorname{Tor}$ szukaj jądra mnożenia przez osiem w $B$, a dla $\operatorname{Ext}$ ilorazu $B/8B$.}
\item Niech $R$ będzie dziedziną z ciałem ułamków $K$. Uzasadnij, że ranga $R^r\oplus T$ jest równa $r$, gdy każdy element $T$ jest zabijany przez jakiś niezerowy element $R$.
\wskazowka{W $K\otimes_R T$ prosty tensor $1\otimes t$ jest zerowy, jeśli $rt=0$ dla $r\ne0$.}
\end{enumerate}

\chapter{Algebra homologiczna: kompleksy i (ko)homologie}
\begin{enumerate}
\item \textbf{Kompleks z trzema stopniami.} Niech $C_2=\Z^2$, $C_1=\Z^3$, $C_0=\Z$, przy czym $\partial_2(a,b)=(2a,0,3b)$, $\partial_1(x,y,z)=y$. Wyznacz wszystkie grupy $H_k(C)$ i wskaż reprezentantów ich generatorów.
\item \textbf{Dobieranie różniczki.} Niech $C_2=\Z$, $C_1=\Z^2$, $C_0=\Z$, $\partial_2(1)=(4,6)$. Wyznacz wszystkie homomorfizmy $\partial_1:\Z^2\to\Z$ spełniające $\partial_1\partial_2=0$. Wybierz taki, który ma obraz $2\Z$, i oblicz homologię.
\wskazowka{Warunek $\partial_1\partial_2=0$ to pojedyncze równanie diofantyczne na współczynniki funkcjonału. Do homologii użyj bazy $\ker\partial_1$.}
\item \textbf{Dwie redukcje współczynników.} Dla kompleksu $C_1=\Z^2\xrightarrow{\operatorname{diag}(4,9)}C_0=\Z^2$ oblicz homologię nad $\Z$, po tensorowaniu z $\Z/2\Z$ oraz po tensorowaniu z $\Z/3\Z$. Wskaż zjawiska, których nie wykrywa samo tensorowanie grup homologii.
\wskazowka{Po tensorowaniu różniczki redukują się modulo $p$. Współczynniki mogą tworzyć nowe cykle; porównaj wynik z członem $\operatorname{Tor}$ w twierdzeniu o współczynnikach uniwersalnych.}
\item \textbf{Kohomologia dualnego kompleksu.} Dla $C_1=\Z^2\xrightarrow{\begin{psmallmatrix}3&6\\0&5\end{psmallmatrix}}C_0=\Z^2$ oblicz $H_*(C)$ i $H^*(\Hom(C,\Z))$. Wyjaśnij, w jakich stopniach pojawia się torsja.
\item \textbf{Łańcuchowa mapa diagonalna.} Niech $C_1=\Z\xrightarrow{\cdot12}C_0=\Z$. Dla liczb całkowitych $a,b$ rozstrzygnij, kiedy $f_1=\cdot a$, $f_0=\cdot b$ wyznaczają endomorfizm kompleksu. Wyznacz wszystkie możliwe endomorfizmy indukowane na $H_0(C)$.
\item \textbf{Stożek macierzy.} Zbuduj stożek odwzorowania $f:\Z^2[0]\to\Z^2[0]$ o macierzy $\begin{psmallmatrix}2&1\\0&3\end{psmallmatrix}$. Oblicz jego homologię i powiąż z jądrem oraz kokernlem $f$.
\wskazowka{Dla kompleksu skoncentrowanego w stopniu zerowym stożek ma tylko dwa niezerowe stopnie. Sprawdź dokładnie przyjętą konwencję znaków różniczki stożka.}
\item \textbf{Acykliczny, lecz nieściągalny kompleks.} Zbuduj dokładny kompleks długości trzy z modułów nad $\Z/4\Z$, który nie jest łańcuchowo homotopijny z zerem. Uzasadnij brak homotopii, korzystając z konsekwencji rozszczepialności.
\wskazowka{Wykorzystaj nierozszczepialny krótki ciąg dokładny $0\to 2R\to R\to R/(2)\to0$ dla $R=\Z/4\Z$. Kompleks ściągalny wymuszałby istnienie odpowiedniego rozszczepienia.}
\item \textbf{Homotopia łańcuchowa w macierzach.} Na kompleksie $0\to\R\xrightarrow{\binom{1}{0}}\R^2\xrightarrow{(0\ \ 1)}\R\to0$ skonstruuj jawny operator stopnia $+1$, $h$, spełniający $\partial h+h\partial=\id$.
\wskazowka{Zacznij od prawego końca kompleksu, dobierając prawą odwrotność surjekcji; następnie zbuduj homotopię na składniku dopełniającym obraz.}
\item \textbf{Pięciowyrazowy ciąg.} Załóż, że $0\to A\to\Z^3\xrightarrow{(2,4,6)}\Z\to B\to0$ jest dokładny. Oblicz $A$ i $B$ i ustal, czy $A$ jest wolny. Wskaż, jakie informacje daje sama znajomość rzędów homomorfizmów.
\item \textbf{Długi ciąg na konkretnym ilorazie.} Niech $B$ będzie kompleksem $B_1=\Z\xrightarrow{\cdot6}B_0=\Z$, a $A\subset B$ podkompleksem zadanym przez $A_1=0$ i $A_0=3\Z$. Oblicz homologie $A$, $B$, $B/A$ oraz homomorfizm łączący.
\wskazowka{Oblicz najpierw $B/A$ w każdym stopniu. Homomorfizm łączący polega na podniesieniu cyklu z ilorazu, zastosowaniu różniczki i odczytaniu klasy w $A$.}
\item \textbf{Filtracja trójstopniowa.} Skonstruuj filtrację kompleksu $C$ z zadania 1 przez podkompleksy, w której w kolejnych krokach dodajesz generatory w stopniach $0,1,2$. Zapisz względne grupy homologii kolejnych par i wyjaśnij, kiedy nowy generator zabija wcześniejszą klasę.
\wskazowka{W filtracji przez szkielety każda nowa komórka daje generator względny, ale jego brzeg może trafić do wcześniejszego filtru.}
\item \textbf{Mayer--Vietoris dla podkompleksów.} Niech $C$ będzie kompleksem symplicjalnym dwóch trójkątów mających jedną wspólną krawędź. Przyjmij za $A,B$ podkompleksy pojedynczych wypełnionych trójkątów, tak że $A\cup B=C$ i $A\cap B$ jest jedną krawędzią. Wypisz łańcuchowy ciąg Mayera--Vietorisa i oblicz jego homomorfizm łączący w stopniu $1$.
\item \textbf{Charakterystyka Eulera bez obliczania homologii.} Kompleks skończonych przestrzeni nad ciałem ma wymiary $(\dim C_3,\dim C_2,\dim C_1,\dim C_0)=(2,5,7,3)$. Załóż, że $H_3=H_0=0$ i $\dim H_2=1$. Wyznacz $\dim H_1$ albo wykaż, że dane są sprzeczne.
\item \textbf{Zastosowanie lematu pięciu.} W przemiennym diagramie długich ciągów dokładnych załóż, że znasz izomorfizmy na czterech spośród pięciu kolejnych grup. Sformułuj precyzyjnie, jakie warunki wystarczą, by uzyskać izomorfizm w środku. Znajdź przykład pokazujący, że osłabienie jednego warunku do samej różnowartościowości nie zawsze wystarcza.
\wskazowka{Rozdziel twierdzenie pięciu na dowody różnowartościowości i surjektywności za pomocą pościgu za diagramem; osłabione założenie sprawdź na prostych ciągach zerowych.}
\item \textbf{Torsja a charakterystyka Eulera.} Porównaj kompleksy $0\to\Z\xrightarrow{\cdot m}\Z\to0$ dla $m=0,1,7$. Oblicz homologię i charakterystykę Eulera każdego z nich. Wytłumacz, dlaczego ostatnia nie odróżnia tych przypadków.
\item \textbf{Uwolnienie torsji przez współczynniki.} Dla kompleksu $0\to\Z\xrightarrow{6}\Z\to0$ oblicz jego homologie po tensorowaniu z $\Z/4$ i $\Z/9$. Porównaj z przewidywaniami twierdzenia o współczynnikach uniwersalnych i odpowiednimi grupami $\operatorname{Tor}$.
\wskazowka{Dla homomorfizmu mnożenia przez $6$ po redukcji modulo $m$ oblicz oddzielnie jądro i obraz. Porównaj z $\operatorname{Tor}_1^{\Z}(\Z/6,\Z/m)$.}
\end{enumerate}


\chapter{Homotopia, nakrycia i CW-kompleksy}
\begin{enumerate}
\item \textbf{Uzupełniony graf.} Niech $X$ powstanie z okręgu przez dołączenie odcinka, którego oba końce przyklejamy do dwóch różnych punktów okręgu. Wyznacz $\pi_1(X)$ oraz liczbę niezależnych cykli grafu.
\item \textbf{Trzy dziury w płaszczyźnie.} Dla $X=\R^2\setminus\{(0,0),(1,0),(0,1)\}$ skonstruuj konkretne pętle generujące $\pi_1(X)$, oblicz obraz komutatora pierwszych dwóch generatorów w abelianizacji i uzasadnij, dlaczego komutator nie jest trywialny w $\pi_1(X)$.
\item \textbf{Przestrzeń z dwoma komórkami.} Niech $X_m=S^1\cup_{z\mapsto z^m}D^2$ dla $m\ge2$. Oblicz $\pi_1(X_m)$. Następnie dołącz drugą $2$-komórkę odwzorowaniem stopnia $n$ i wyznacz grupę podstawową nowej przestrzeni.
\item \textbf{Nakrycie skończonego grafu.} Zbuduj spójne nakrycie stopnia $3$ bukietu dwóch okręgów, opisując działanie generatorów na trzy punkty włókna przez permutacje. Wyznacz grupę automorfizmów dla swojego przykładu.
\wskazowka{Nakrycie grafu zakoduj dwiema permutacjami włókna. Grupa przekształceń pokładowych jest centralizatorem działania monodromii w $S_3$.}
\item \textbf{Nakrycia torusa z macierzy.} Niech $A=\begin{psmallmatrix}2&1\\0&3\end{psmallmatrix}$, $f_A:\T^2\to\T^2$, $f_A([x])=[Ax]$. Wyznacz stopień nakrycia, grupę przekształceń pokładowych i podgrupę $\pi_1(\T^2)$ odpowiadającą nakryciu.
\wskazowka{Podgrupą odpowiadającą nakryciu jest $A\Z^2\subset\Z^2$. Grupę przekształceń pokładowych znajdziesz jako $\Z^2/A\Z^2$.}
\item \textbf{Podnoszenie do nakrycia torusa.} Dla macierzy całkowitych $A,B\in M_2(\Z)$ i $\det A\ne0$ znajdź warunek konieczny i dostateczny na istnienie $g:\T^2\to\T^2$ z $f_A\circ g=f_B$. Sprawdź go dla $A=\operatorname{diag}(2,3)$ oraz $B=\begin{psmallmatrix}4&2\\0&6\end{psmallmatrix}$.
\wskazowka{Po przejściu na grupy podstawowe warunek podniesienia sprowadza się do zawierania $B\Z^2\subseteq A\Z^2$; następnie uwzględnij ewentualne przesunięcie podniesienia.}
\item \textbf{Nakrycie niegaloisowskie.} Skonstruuj spójne trójkrotne nakrycie bukietu dwóch okręgów, którego grupa przekształceń pokładowych jest trywialna. Użyj dwóch odpowiednio dobranych permutacji z $S_3$.
\wskazowka{Wybierz permutacje generujące przechodnie działanie z trywialnym centralizatorem, na przykład transpozycję i cykl długości trzy.}
\item \textbf{Prezentacja z komórek.} Zbuduj $2$-kompleks CW o jednej $0$-komórce, dwóch $1$-komórkach i jednej $2$-komórce, której odwzorowanie dołączające opisuje słowo $a^2b^3$. Zapisz prezentację $\pi_1$ i oblicz jego abelianizację.
\wskazowka{Relację dołączającej komórki odczytaj z wyrazu w grupie wolnej. Abelianizację otrzymasz, zastępując ją wektorem wykładników.}
\item \textbf{Połączenie dwóch torusów punktem.} Wyznacz $\pi_1(\T^2\vee\T^2)$ przez prezentację. Rozstrzygnij, czy grupa jest abelowa, i znajdź jej epimorfizm na wolną grupę rzędu dwa.
\item \textbf{Sklejanie walców.} Dwa cylindry $S^1\times[0,1]$ sklej wzdłuż jednego okręgu brzegowego odwzorowaniem stopnia $\pm1$. Ustal typ homotopii wyniku. Co zmieni się, jeśli użyjesz stopnia $2$ i zamiast homeomorfizmu rozważysz iloraz sklejony wzdłuż takiego odwzorowania?
\item \textbf{Wolne grupy i drzewa.} Niech $G$ będzie grafem powstałym z pięciokąta przez dodanie dwóch przekątnych nieprzecinających się poza wierzchołkami. Znajdź drzewo rozpinające i bazę $\pi_1(G)$; opisz zmianę bazy przy wyborze innego drzewa.
\item \textbf{Pętla a stopień.} Dla $f:S^1\to S^1$, $f(e^{it})=e^{i(5t+\sin t)}$, oblicz $f_*:\pi_1(S^1)\to\pi_1(S^1)$. Rozstrzygnij, przez które z nakryć $z\mapsto z^n$ odwzorowanie $f$ ma podniesienie.
\item \textbf{Grupa podstawowa sumy spójnej powierzchni.} Skonstruuj CW-model orientowalnej powierzchni rodzaju $2$ z jednym wierzchołkiem i jedną $2$-komórką. Wyznacz prezentację grupy podstawowej i jej abelianizację.
\wskazowka{Użyj standardowego wielokąta ośmiobocznego z parami utożsamianych boków; słowo brzegu daje relację komórki dwuwymiarowej.}
\item \textbf{Graf a przyczepienie komórki.} Zacznij od bukietu trzech okręgów z generatorami $a,b,c$. Dołącz $2$-komórki o słowach $ab^{-1}$ i $bc^{-1}$. Oblicz grupę podstawową i porównaj ją z wynikiem po dołączeniu dodatkowej komórki o słowie $a^4$.
\item \textbf{Kontrprzykład do intuicji o $\pi_1$.} Znajdź dwa skończone CW-kompleksy o izomorficznej grupie podstawowej, które nie są homotopijnie równoważne. Wykaż brak równoważności za pomocą odpowiedniego niezmiennika z dalszych rozdziałów.
\item \textbf{Nakrycia butelki Kleina.} Dla prezentacji $\pi_1(K)=\langle a,b\mid aba^{-1}=b^{-1}\rangle$ wyznacz jądro homomorfizmu $\pi_1(K)\to\Z/2$ danego przez $a\mapsto1$, $b\mapsto0$. Zidentyfikuj odpowiadające mu dwukrotne nakrycie i oblicz jego grupę podstawową.
\wskazowka{Rozważ homomorfizm określający orientację; w jądrze sprawdź relację pomiędzy $a^2$ oraz $b$.}
\end{enumerate}


\chapter{Homologia przestrzeni}
\begin{enumerate}
\item \textbf{Pięciokąt bez wypełnienia.} Zbuduj macierz różniczki dla kompleksu symplicjalnego będącego brzegiem pięciokąta. Oblicz $H_0,H_1$ nad $\Z$ i porównaj wyniki po dodaniu wszystkich trójkątów triangulujących wnętrze.
\item \textbf{Dwa okręgi połączone krawędzią.} Dla grafu utworzonego z dwóch rozłącznych trójkątnych cykli przez dołączenie jednej krawędzi łączącej ich wierzchołki oblicz homologie symplicjalne i podaj jawny układ generatorów $H_1$.
\item \textbf{Kompleks z torsją.} Niech $X_m=S^1\cup_m e^2$ dla $m\ge2$. Zapisz wszystkie komórkowe operatory brzegowe i oblicz $H_k(X_m;\Z)$ oraz $H_k(X_m;\Z/p\Z)$ dla liczby pierwszej $p$.
\wskazowka{Użyj kompleksu komórkowego z pojedynczą różniczką mnożenia przez $m$. Po redukcji modulo $p$ operator może stać się zerowy.}
\item \textbf{Dwa dołączenia do okręgu.} Dla $X=S^1\cup_m e^2\cup_n e^2$ oblicz $H_2(X;\Z)$ i $H_1(X;\Z)$ jako funkcje $m,n\in\Z$. Wyjaśnij, skąd bierze się wolny składnik w $H_2$.
\wskazowka{Operator $\partial_2$ jest wierszem $(m,n)$. Jego jądro i kokernel oblicz przez największy wspólny dzielnik.}
\item \textbf{Kompleks o zadanej macierzy.} Skonstruuj $2$-wymiarowy CW-kompleks z jedną $0$-komórką i dwiema $1$- oraz dwiema $2$-komórkami, którego różniczka $\partial_2$ ma macierz $\begin{psmallmatrix}2&4\\0&6\end{psmallmatrix}$. Oblicz homologie całkowite.
\wskazowka{Macierz różniczki sprowadź do postaci Smitha operacjami elementarnymi nad $\Z$; odpowiednie słowa dołączania można wybrać w wolnej grupie dwóch generatorów.}
\item \textbf{Homologia powierzchni z brzegiem.} Rozważ orientowalną powierzchnię rodzaju $g$ z $b\ge1$ okręgami brzegowymi. Wyznacz jej grupy homologii i charakterystykę Eulera za pomocą odpowiedniego modelu CW.
\item \textbf{Produkt z okręgiem.} Oblicz $H_*(S^2\times S^1;\Z)$ na podstawie struktury CW produktu lub wzoru Künnetha. Wskaż konkretne cykle reprezentujące niezerowe klasy.
\item \textbf{Klein i współczynniki.} Oblicz $H_*(K;\Z/2\Z)$ dla butelki Kleina, zaczynając od jej komórkowego kompleksu nad $\Z$. Porównaj z $H_*(K;\Q)$ i wyjaśnij różnicę.
\wskazowka{W standardowym modelu CW butelki Kleina jedyna nietrywialna różniczka jest mnożeniem przez $2$. Zredukuj ją modulo wybrany pierścień współczynników.}
\item \textbf{Homologia po sklejeniu dwóch dysków.} Sklej dwa dyski $D^2$ wzdłuż całych brzegów przez odwzorowanie stopnia $-1$. Zidentyfikuj przestrzeń ilorazową i oblicz jej homologię bez korzystania z gotowej nazwy tej przestrzeni.
\item \textbf{Zawieszenie.} Dla spójnego CW-kompleksu $X$ oblicz $\widetilde H_k(\Sigma X)$ przez Mayera--Vietorisa. Zastosuj do $X=S^1\vee S^2$, podając wymiary wszystkich grup.
\wskazowka{Pokryj zawieszenie dwoma stożkami, których przecięcie deformacyjnie retraktuje się na $X$. Użyj zredukowanego ciągu Mayera--Vietorisa.}
\item \textbf{Para walca.} Dla $M=S^1\times[0,1]$ i $A=S^1\times\{0,1\}$ oblicz $H_*(M,A;\Z)$ i opisz homomorfizm łączący do $H_0(A)$.
\item \textbf{Homologia z identyfikacją brzegów.} Przestrzeń $X$ otrzymuje się z kwadratu przez sklejenie pionowych boków zgodnie z orientacją, a poziomych z odwróceniem orientacji. Zapisz komórkowy kompleks i oblicz $H_*(X;\Z/3\Z)$.
\item \textbf{Homologia lokalna w węźle grafu.} Niech $X$ będzie grafem z wierzchołkiem $v$ o walencji $5$. Oblicz $H_k(X,X\setminus\{v\};\Z)$ i porównaj z lokalną homologią punktu wewnętrznego krawędzi.
\wskazowka{Wytnij małe gwiazdowe otoczenie wierzchołka. Homologia lokalna w stopniu jeden odpowiada zredukowanej homologii jego łącza.}
\item \textbf{Stopień i stożek odwzorowania.} Niech $f:S^2\to S^2$ ma stopień $m\ne0$. Dołącz do $S^2$ jedną $3$-komórkę przez $f$. Oblicz wszystkie grupy homologii powstałej przestrzeni i sprawdź zależność od $m$.
\wskazowka{W komórkowym kompleksie różniczka od nowej $3$-komórki do $2$-komórki jest mnożeniem przez stopień $m$.}
\item \textbf{Przestrzeń z dwoma torusami.} Rozważ $X=(T^2\sqcup T^2)/(p\sim q)$, gdzie identyfikujesz jeden punkt w każdym torusie. Oblicz homologie $X$ i porównaj je z homologiami rozłącznej sumy dwóch torusów.
\wskazowka{Identyfikacja punktów daje sumę klinową torusów; skorzystaj ze zredukowanych homologii sumy klinowej.}
\end{enumerate}


\chapter{Wyższe grupy homotopii i związek z homologią}
\begin{enumerate}
\item \textbf{Wyższe grupy produktu.} Dla $n\ge2$ wykaż, że $\pi_n(X\times Y)\cong\pi_n(X)\oplus\pi_n(Y)$ dla przestrzeni bazowanych. Zastosuj do $S^2\times S^3$, wyznaczając $\pi_2$ oraz $\pi_3$; w drugim przypadku wolno przyjąć klasyczny wynik $\pi_3(S^2)\cong\Z$.
\item \textbf{Nakrycie i $\pi_2$.} Oblicz $\pi_2(\mathbb{RP}^3)$ za pomocą nakrycia $S^3\to\mathbb{RP}^3$. Porównaj wynik z $H_2(\mathbb{RP}^3;\Z)$.
\wskazowka{Nakrycie indukuje izomorfizmy grup homotopii w stopniach co najmniej dwa; oblicz je po przejściu do nakrycia uniwersalnego.}
\item \textbf{Przestrzeń Moore'a.} Dla $m\ge2$ rozważ $M_m=S^2\cup_m e^3$. Oblicz $\pi_2(M_m)$ za pomocą twierdzenia Hurewicza oraz $H_2(M_m;\Z)$ z kompleksu komórkowego.
\wskazowka{Przestrzeń Moore'a buduje się przez dołączenie komórki z odwzorowaniem określonego stopnia. Zapisz odpowiadający kompleks komórkowy i użyj twierdzenia Hurewicza dopiero po sprawdzeniu spójności.}
\item \textbf{Wykrywanie nietrywialnej wyższej grupy.} Niech $X=S^2\vee S^3$. Oblicz $\pi_2(X)$, $H_2(X)$ i wskaż klasę, którą wykrywa odwzorowanie Hurewicza. Uzasadnij, dlaczego samo $H_2$ nie wyznacza wszystkich wyższych grup homotopii.
\item \textbf{Przyczepienie komórki w stopniu wyższym.} Niech $X=S^3\cup_f e^4$ dla $f:S^3\to S^3$ o stopniu $7$. Oblicz $\pi_3(X)$ i $H_3(X;\Z)$, wyjaśniając, jak działa dołączenie komórki na generator sfery.
\wskazowka{Dołączenie komórki $n$-wymiarowej nie zmienia szkieletu niższego wymiaru; rozważ długi ciąg homotopii pary, zamiast tylko homologii.}
\item \textbf{Długi ciąg homotopii produktu z parą.} Dla pary $(S^2\times D^3,S^2\times S^2)$ oblicz $\pi_3$ względną, wykorzystując równoważność homotopijną i długi ciąg pary. Opisz obraz generatora w $\pi_2(S^2\times S^2)$.
\item \textbf{Stopień i $\pi_n$.} Niech $f:S^n\to S^n$, $n\ge2$, ma stopień $-3$. Oblicz homomorfizm $f_*:\pi_n(S^n)\to\pi_n(S^n)$ i rozstrzygnij, czy $f$ może być równoważnością homotopijną.
\item \textbf{Przyczepienie komórki a spójność.} Niech $X$ będzie jednospójny, a $Y=X\cup_f D^{n+1}$, $n\ge3$. Wyznacz zakres stopni $k$, w których $\pi_k(X)\to\pi_k(Y)$ musi być izomorfizmem, oraz stopień, w którym może nie być surjekcją lub injekcją. Uzasadnij w oparciu o względną aproksymację komórkową.
\wskazowka{Względne grupy homotopii pary $(X\cup_f D^n,X)$ są kontrolowane przez wymiar nowej komórki. Wykorzystaj wynik o spójności pary CW.}
\item \textbf{Przestrzeń o zerowych grupach niższego rzędu.} Rozważ CW-kompleks $X=S^4\vee S^4\vee S^5$. Oblicz $\pi_i(X)$ dla $i\le4$ oraz homomorfizm Hurewicza w stopniu $4$.
\item \textbf{Abelianizacja i odwzorowanie do okręgu.} Niech $X$ będzie powierzchnią orientowalną rodzaju $2$. Wyznacz $\pi_1(X)_{\mathrm{ab}}$ i opisz możliwe homomorfizmy $\pi_1(X)\to\Z$ jako obrazy indukowane przez odwzorowania $X\to S^1$.
\item \textbf{Porównanie dwóch sferycznych nakryć.} Oblicz $\pi_3(\mathbb{RP}^2)$ i $\pi_3(\mathbb{RP}^3)$, odwołując się do odpowiednich nakryć i przyjmując klasyczny wynik Hopfa $\pi_3(S^2)\cong\Z$. Wyjaśnij, dlaczego ostatniego wyniku nie można uzyskać przez bezpośrednie zastosowanie twierdzenia Hurewicza do $S^2$ w stopniu $3$.
\item \textbf{Działanie $\pi_1$ w przestrzeni rzutowej.} Dla parzystego $n\ge2$ określ działanie $\pi_1(\mathbb{RP}^n)$ na $\pi_n(\mathbb{RP}^n)\cong\Z$. Porównaj przypadek nieparzystego $n$ i uzasadnij rolę stopnia odwzorowania antypodycznego.
\wskazowka{Działanie grupy podstawowej na $\pi_2$ można odczytać z działania transformacji pokładowych nakrycia uniwersalnego.}
\item \textbf{Zawieszenie i pierwsza homotopia.} Dla $X=S^1\vee S^1$ oblicz grupy $\pi_2(\Sigma X)$ i $H_2(\Sigma X)$. Wskaż, w którym miejscu stosujesz jednospójność zawieszenia.
\item \textbf{Homologia nie wykrywa grupy podstawowej.} Porównaj $X=S^1\vee S^2$ i $Y=S^1\times S^2$ pod kątem $H_1,H_2,H_3$ oraz $\pi_1$. Znajdź niezmiennik wykrywający różnicę między ich typami homotopii.
\item \textbf{Homologia wymierna to za mało.} Niech $f:\T^2\to\T^2$ będzie odwzorowaniem indukowanym przez macierz $\operatorname{diag}(2,3)$. Oblicz $f_*$ na homologii całkowitej i wymiernej we wszystkich stopniach. Wyjaśnij, dlaczego izomorfizmy w homologii wymiernej nie pozwalają zastosować homologicznego kryterium Whiteheada.
\wskazowka{Homologia wymierna pomija torsję. Znajdź przykład komórkowy, w którym różniczka jest mnożeniem przez liczbę niezerową, lecz nieodwracalną nad $\Z$.}
\item \textbf{Włóknienie Hopfa jako narzędzie obliczeń.} Korzystając z długiego ciągu dokładnego homotopii włóknienia $S^1\hookrightarrow S^3\to S^2$, wyznacz $\pi_2(S^2)$ i $\pi_3(S^2)$ z uwzględnieniem grup homotopii sfer występujących w ciągu. Wskaż homomorfizm brzegowy i jego interpretację.
\wskazowka{Zapisz długi ciąg homotopii włóknienia $S^1\to S^3\to S^2$ w okolicy badanej grupy, pamiętając o znikaniu wyższych grup homotopii $S^1$.}
\end{enumerate}


\chapter{Kohomologia singularna i iloczyn kubkowy}
\begin{enumerate}
\item Z definicji kołańcuchów udowodnij $H^0(X;G)\cong\prod_{C\in\pi_0(X)}G$, gdzie $\pi_0(X)$ oznacza zbiór składowych łukowych.
\item Wyznacz $H^*(S^n;\Z)$ i porównaj wyniki z $H_n(S^n;\Z)$ dla $n\geq1$.
\item Oblicz $H^*(\R P^2;\Z)$, stosując twierdzenie o współczynnikach uniwersalnych, a następnie $H^*(\R P^2;\Z/2\Z)$.
\wskazowka{W kompleksie komórkowym $\R P^2$ różniczka stopnia dwa to mnożenie przez $2$.}
\item Wyprowadź z reguły Leibniza fakt, że iloczyn kubkowy jest dobrze określony na klasach kohomologii.
\item Oblicz pierścień $H^*(S^1\times S^1;\Z)$ oraz zidentyfikuj generator drugiej kohomologii.
\item Wyjaśnij, dlaczego sama znajomość grup $H^n(X;\Z)$ nie wystarcza do ustalenia pierścienia $H^*(X;\Z)$.
\item Oblicz $\operatorname{Ext}^1_{\Z}(\Z/m\Z,\Z/n\Z)$ bez cytowania gotowego wzoru.
\wskazowka{Zastosuj $\operatorname{Hom}(-,\Z/n\Z)$ do $0\to\Z\xrightarrow{m}\Z\to\Z/m\Z\to0$.}
\item Użyj wzoru Künnetha do obliczenia grup homologii $S^2\times S^3$.
\end{enumerate}
\chapter{Formy różniczkowe, orientacja i kohomologia de Rhama}
\begin{enumerate}
\item \textbf{Zamkniętość formy z parametrem.} Dla $\omega=(ax^2y+y)\,dx+(x^3+bx)\,dy$ na $\R^2$ wyznacz parametry $a,b\in\R$, dla których $\omega$ jest zamknięta. Znajdź funkcję pierwotną w tych przypadkach.
\wskazowka{Wypisz $d\alpha$ i osobno zbierz współczynniki każdego niezależnego dwuwektora bazowego. Otrzymany warunek będzie równaniem na parametr.}
\item \textbf{Cofnięcie formy objętości.} Niech $F(u,v,w)=(u+v,uv,w+u^2)$. Oblicz $F^*(dx\wedge dy\wedge dz)$ i znajdź zbiór, w którym $F$ zachowuje orientację, odwraca ją lub nie jest lokalnym dyfeomorfizmem.
\item \textbf{Tożsamość Cartana w przykładzie.} Dla $X=y\partial_x-x\partial_y$ i $\omega=x\,dy-y\,dx$ na $\R^2$ oblicz $i_X\omega$, $i_X\dd\omega$ i $\mathcal L_X\omega$. Porównaj wynik z zachowaniem $\omega$ przez obroty.
\item \textbf{Forma na pierścieniu.} Dla $A=\{1<x^2+y^2<9\}$ i $\omega=((2x-y)\,dx+(x+2y)\,dy)/(x^2+y^2)$ rozłóż $\omega$ na część dokładną i wielokrotność formy kąta. Oblicz okres po dodatnio zorientowanym okręgu $r=2$.
\wskazowka{Porównaj całkę po okręgu otaczającym dziurę z całką po pętli ściągalnej. Nie wystarcza jedynie sprawdzenie $d\alpha=0$.}
\item \textbf{Dwie dziury, dwa okresy.} Na $U=\R^2\setminus\{(0,0),(2,0)\}$ zdefiniuj formy $\omega_1,\omega_2$ jako formy kąta wokół dwóch usuniętych punktów. Udowodnij liniową niezależność ich klas w $H^1_{\mathrm{dR}}(U)$, obliczając całki po małych okręgach.
\wskazowka{Wybierz bazę $H_1$ reprezentowaną przez małe dodatnio zorientowane okręgi wokół obu dziur. Okresy są współrzędnymi klasy de Rhama.}
\item \textbf{Stokes na pierścieniu.} Dla $\alpha=(x^3-y)\,dy$ oblicz $\int_{\partial A}\alpha$ oraz $\int_A\dd\alpha$, gdzie $A=\{1\le x^2+y^2\le4\}$. Ustal poprawną orientację wewnętrznej składowej brzegu.
\item \textbf{Całkowanie na walcu.} Rozważ $M=S^1\times[0,2]$ z orientacją $d\theta\wedge dt$ i formę $\omega=t\sin\theta\,d\theta+\cos\theta\,dt$. Oblicz $\dd\omega$ i obie strony twierdzenia Stokesa.
\item \textbf{Forma powierzchni na torusie.} Dla płaskiego torusa $\T^2=\R^2/\Z^2$ sprawdź, czy forma $\omega=(2+\sin(2\pi x))\,dx\wedge dy$ jest dokładna. Oblicz jej klasę w $H^2_{\mathrm{dR}}(\T^2)$ względem generatora $[dx\wedge dy]$.
\item \textbf{Okresy form w stopniu jeden.} Na $\T^2$ rozważ $\alpha=a\,dx+b\,dy+\dd f$, gdzie $f$ jest gładka i okresowa. Oblicz całki po dwóch podstawowych pętlach i rozstrzygnij, kiedy $\alpha$ jest dokładna.
\wskazowka{Sprowadź badanie dokładności do okresów po generatorach pierwszej homologii; do zbudowania potencjału użyj całki po drodze od punktu bazowego.}
\item \textbf{Kohomologia cylindra z usuniętą linią.} Oblicz $H^k_{\mathrm{dR}}((S^1\setminus\{p\})\times\R)$ i porównaj z kohomologią $S^1\times\R$. Zbuduj odpowiednie retrakcje deformacyjne.
\wskazowka{Najpierw zidentyfikuj składowe przestrzeni po usunięciu wskazanego podzbioru; kohomologię liczy się osobno na każdej składowej.}
\item \textbf{Forma magnetyczna na $\R^3$.} Dla $\omega=x\,dy\wedge dz+y\,dz\wedge dx+z\,dx\wedge dy$ oblicz $\dd\omega$ i znajdź $2$-formę $\eta$ o zadanej pochodnej $\dd\eta=\dd\omega$ i prostszych współczynnikach. Zinterpretuj niejednoznaczność wyboru $\eta$.
\item \textbf{Orientacja regularnej poziomicy.} Niech $M=\{(x,y,z):x^2+2y^2+3z^2=1\}$. Dla orientacji otaczającej $\R^3$ i normalnej zewnętrznej wyznacz dodatnią orientację $T_pM$ w punkcie $p=(1,0,0)$ i porównaj ją z orientacją parametryzacji lokalnej.
\wskazowka{Orientację regularnej poziomicy można opisać przez warunek, by dodatni wektor normalny i dodatnia baza przestrzeni stycznej dawały orientację otoczenia.}
\item \textbf{Zależność od orientacji.} Na torusie $\T^2$ rozważ $\omega=\cos(2\pi x)\,dx\wedge dy$ i $\eta=(1+\cos(2\pi x))\,dx\wedge dy$. Oblicz całki dla obu orientacji, a następnie rozstrzygnij, które z tych form mogą być dokładne.
\item \textbf{Kohomologia pary spodni.} Niech $P$ będzie powierzchnią powstałą ze sfery przez usunięcie wnętrz trzech rozłącznych domkniętych dysków. Korzystając z ciągu Mayera--Vietorisa, oblicz $H^k_{\mathrm{dR}}(P)$ i znajdź relację między okresami zamkniętej $1$-formy po trzech zorientowanych składowych brzegu.
\wskazowka{Para spodni retraktuje się na graf z dwiema niezależnymi pętlami. Uwzględnij relację między trzema krzywymi brzegowymi przy całkowaniu zamkniętych form.}
\item \textbf{Dwie osobliwości formy powierzchniowej.} Dla $p=(0,0,1)$ i $q=(0,0,-1)$ niech $U=\R^3\setminus\{p,q\}$ oraz $\omega_r=\bigl((x-r_1)\,dy\wedge dz+(y-r_2)\,dz\wedge dx+(z-r_3)\,dx\wedge dy\bigr)/\|x-r\|^3$ dla $r=p,q$. Oblicz okresy formy $2\omega_p-\omega_q$ po małych dodatnio zorientowanych sferach wokół $p$ i $q$. Wyznacz wymiar podprzestrzeni $H^2_{\mathrm{dR}}(U)$ rozpiętej przez jej klasy składowe.
\wskazowka{Oddziel lokalne okresy wokół poszczególnych osobliwości od okresu po dużej pętli obejmującej je obie; zastosuj addytywność klas homologii.}
\item \textbf{Stopień z całki formy objętości.} Na torusie $T^2=\R^2/\Z^2$ rozważ odwzorowanie $f([x,y])=([2x+y],[x+3y])$. Oblicz $f^*(\dd x\wedge\dd y)$ i wyznacz stopień $f$ przez porównanie całek formy objętości. Wyjaśnij, dlaczego wynik zgadza się z liczbą arkuszy nakrycia i orientacją.
\wskazowka{Do obliczenia stopnia użyj $\int_M f^*\omega=\deg(f)\int_N\omega$ dla znormalizowanej formy objętości; wyjaśnij wpływ orientacji.}
\end{enumerate}


\chapter{Przecięcia, dualność Poincarégo i klasy charakterystyczne}
\begin{enumerate}
\item \textbf{Iloczyn kubkowy na produkcie.} Na $M=S^1\times S^2$ oznacz przez $a\in H^1(M;\Z)$ i $b\in H^2(M;\Z)$ klasy ściągnięte z czynników. Oblicz $a\smile b$, $a^2$, $b^2$ i opisz strukturę pierścienia $H^*(M;\Z)$.
\item \textbf{Pierścień kohomologii torusa trójwymiarowego.} Niech $x,y,z$ będą standardowymi generatorami $H^1(T^3;\Z)$. Wyznacz bazę $H^2$ i $H^3$ z ich iloczynów kubkowych. Oblicz $(x+2y)\smile(y+3z)\smile(z+x)$.
\wskazowka{Pierścień kohomologii $T^3$ jest algebrą zewnętrzną na trzech generatorach stopnia jeden. Przed mnożeniem uprość wyrażenie przez antyprzemienność.}
\item \textbf{Krzywe na torusie.} Dwie zorientowane krzywe proste na $T^2$ reprezentują klasy $(2,1)$ i $(1,3)$ w $H_1(T^2;\Z)$. Oblicz ich algebraiczną liczbę przecięcia i minimalną możliwą liczbę przecięć poprzecznych.
\item \textbf{Dualność na $S^1\times S^2$.} Wyznacz obrazy generatorów $H^1$ i $H^2$ przez dualność Poincarégo $H^k(M;\Z)\to H_{3-k}(M;\Z)$. Przedstaw je jako konkretne podrozmaitości.
\item \textbf{Forma przecięcia produktu.} Dla $M=S^2\times S^2$ wybierz klasy $A=[S^2\times\{p\}]$, $B=[\{q\}\times S^2]$. Oblicz macierz formy przecięcia na $H_2(M;\Z)$ i samoprzecięcia klas $A+B$ oraz $2A-3B$.
\item \textbf{Brak orientacji i homologia najwyższa.} Dla butelki Kleina oblicz $H_2(K;\Z)$ oraz $H_2(K;\Z/2\Z)$. Wyjaśnij, jakie podstawowe klasy fundamentalne istnieją i jak zmienia się sformułowanie dualności Poincarégo.
\item \textbf{Funkcja sklejająca wiązki.} Nad $S^2$ sklej dwie trywialne rzeczywiste wiązki rzędu dwa za pomocą $g_k:S^1\to SO(2)$, $g_k(e^{it})=R_{kt}$. Dla jakich $k$ otrzymujesz orientowalną wiązkę? Oblicz jej liczbę Eulera zgodnie z ustaloną orientacją.
\wskazowka{Wiązki orientowane płaszczyzn nad $S^2$ klasyfikuje stopień funkcji sklejającej $S^1\to SO(2)$. Liczba Eulera mierzy ten stopień z ustaloną konwencją znaku.}
\item \textbf{Klasa Eulera jako przeszkoda.} Dla zorientowanej wiązki rzędu dwa $E_k\to S^2$ z zadania 7 rozstrzygnij, kiedy istnieje przekrój nigdzie niezerowy. Porównaj odpowiedź z możliwością rozszczepienia $E_k\cong\underline\R\oplus L$.
\item \textbf{Samoprzecięcie przekroju zerowego.} Niech $E_k\to S^2$ będzie wiązką z zadania 7. Wyznacz liczbę samoprzecięcia sekcji zerowej we wnętrzu przestrzeni całkowitej $E_k$, przy przyjętej orientacji.
\item \textbf{Pierwsza klasa Stiefela--Whitneya.} Na $S^1\times S^1$ rozważ sumę $L_1\oplus L_2$, gdzie $L_i$ jest cofnięciem wiązki Möbiusa przez $i$-tą projekcję. Oblicz $w_1$ i $w_2$ tej wiązki oraz ustal jej orientowalność.
\wskazowka{Użyj wzoru Whitneya na $w_1,w_2$ sumy oraz faktu, że każda cofnięta wiązka Möbiusa ma niezerową klasę $w_1$.}
\item \textbf{Wiązka styczna do produktu sfer.} Korzystając ze wzorów na klasy charakterystyczne sum prostych wyznacz $w_2(T(S^2\times S^2))$. Rozstrzygnij, czy ta rozmaitość dopuszcza strukturę spinową.
\item \textbf{Klasy Cherna sumy prostych.} Niech $L,H$ będą zespolonymi wiązkami liniowymi nad $\mathbb{CP}^2$ z $c_1(L)=h$, $c_1(H)=2h$. Oblicz $c_1$ i $c_2$ wiązki $L\oplus H$ i sprawdź, czy może ona posiadać wszędzie niezerowy przekrój.
\wskazowka{Zastosuj $c(L\oplus H)=(1+c_1(L))(1+c_1(H))$. Niezerowa najwyższa klasa Cherna utrudnia istnienie globalnego nigdzie niezerowego przekroju.}
\item \textbf{Poincaré--Hopf z brzegiem.} Na domkniętym dysku w $\R^2$ rozważ pola $X_1=(x,y)$ i $X_2=(x,-y)$. Oblicz lokalne indeksy ich zer. Zbadaj zachowanie na brzegu i wyjaśnij, dla którego pola można bezpośrednio użyć wersji twierdzenia z polem skierowanym na zewnątrz.
\item \textbf{Indeks gradientu nietypowej funkcji na sferze.} Na standardowej $S^2\subset\R^3$ rozważ funkcję $f(x,y,z)=xy$ i pole $X=\operatorname{grad}_{S^2}f$. Wyznacz jego zera, oblicz indeks każdego z nich z hesjanu $f|_{S^2}$ i porównaj sumę z charakterystyką Eulera.
\wskazowka{Zera gradientu na sferze spełniają warunek, że gradient euklidesowy funkcji jest równoległy do promienia. Indeks wynika ze znaku wyznacznika hesjanu ograniczonego.}
\item \textbf{Forma przecięcia i sygnatura.} Niech $M$ będzie zamkniętą orientowalną $4$-rozmaitością, której forma przecięcia w pewnej bazie ma macierz $\begin{psmallmatrix}2&1\\1&-2\end{psmallmatrix}$. Oblicz wyznacznik, parzystość i sygnaturę formy. Rozstrzygnij, czy jest unimodularna i co mówi to o możliwości realizacji jej jako pełnej formy na wolnej części $H_2(M;\Z)$.
\wskazowka{Forma przecięcia zamkniętej orientowalnej $4$-rozmaitości jest unimodularna na części wolnej drugiej homologii. Sprawdź wyznacznik zanim interpretujesz macierz geometrycznie.}
\item \textbf{Parzystość i klasa Stiefela--Whitneya.} Niech $X$ będzie zamkniętą, zorientowaną, jednospójną $4$-rozmaitością, dla której $Q_X=\operatorname{diag}(1,-1)$. Korzystając ze wzoru Wu, rozstrzygnij, czy $X$ jest spinowa. Porównaj z przypadkiem $Q_X=H$.
\wskazowka{Zastosuj charakterystykę $w_2$ przez kongruencję $Q(x,x)\equiv\langle w_2,x\rangle\pmod2$. Porównaj parzystość obu form.}
\end{enumerate}


\chapter{Fibracje: podnoszenie i obliczenia}
\begin{enumerate}
\item \textbf{Transport.} Dla fibracji $B\times F\to B$ podaj dwa różne podniesienia tej samej drogi zaczynające się w tym samym punkcie. Dlaczego nie przeczy to twierdzeniu o monodromii?
\wskazowka{Przemieszczaj drugi współczynnik wzdłuż dowolnej pętli we włóknie.}
\item \textbf{Nakrycie.} Dla $z\mapsto z^m$ na $S^1$ wyznacz monodromię na zbiorze $m$-tych pierwiastków z jedności oraz mapę $\pi_1(S^1)\to\pi_0(F)$.
\item \textbf{Hopf.} Sprawdź bezpośrednio, że $h(\lambda z,\lambda w)=h(z,w)$ dla $h(z,w)=(2z\overline w,|z|^2-|w|^2)$. Wskaż, dlaczego zamiana $\overline w$ na $w$ psuje ten argument.
\item \textbf{Grupy.} Korzystając z $S^1\to S^3\to S^2$, oblicz $\pi_2(S^2)$ i $\pi_3(S^2)$. Wyjaśnij, skąd pochodzą potrzebne grupy sfery $S^3$.
\item \textbf{Włókno homotopijne.} Oblicz $\operatorname{hofib}(f)$ dla stałej mapy $f:X\to B$ przy ustalonym punkcie bazy, zakładając $X,B$ łukowo spójne.
\wskazowka{Rozdziel współrzędne $x$ i drogę rozpoczynającą się w wartości stałej mapy.}

\item \textbf{Podnoszenie w fibracji trywialnej.} Dla $p:B\times F\to B$ skonstruuj jawnie podniesienie dowolnej homotopii $H:Z\times I\to B$, przy danym podniesieniu czasu zerowego $h:Z\to B\times F$.
\item \textbf{Homomorfizm łączący fibracji Hopfa.} Z długiego ciągu dla $S^1\hookrightarrow S^3\to S^2$ wyznacz dokładnie $\partial:\pi_2(S^2)\to\pi_1(S^1)$. Wyjaśnij, dlaczego jest izomorfizmem.
\wskazowka{Wykorzystaj $\pi_2(S^3)=\pi_1(S^3)=0$ i zachowaj kolejność strzałek.}
\item \textbf{Transport a jednoznaczność.} Wyjaśnij, dlaczego w fibracji Hurewicza transport wzdłuż drogi nie musi być jednoznaczny, choć jego klasa homotopii jest określona. Porównaj to z podnoszeniem drogi w nakryciu.

\end{enumerate}

\chapter{Ciągi spektralne i fibracje Serre'a}
\begin{enumerate}
\item Zdefiniuj $E^0$ oraz $E^1$ dla filtrowanego kompleksu $C_2=\Z x\xrightarrow{\partial}\Z y=C_1$ z $\partial x=3y$, $F_0C=\Z y$, $F_1C=F_0C$, $F_2C=C$. Oblicz wszystkie niezerowe $d^r$ i $E^\infty$.
\wskazowka{Jedyna możliwa różniczka jest $d^2_{2,0}$ i działa przez mnożenie przez trzy.}
\item Wykorzystując definicje $Z^r_{p,q}$ i $B^r_{p,q}$, sprawdź samodzielnie, że $d^r$ nie zależy od wyboru reprezentanta i ma kwadrat zero.
\item Skonstruuj dwie filtrowane grupy abelowe o tej samej stowarzyszonej grupie stopniowanej $\Z/2\oplus\Z/2$, ale nieizomorficznych grupach całkowitych.
\item Dla $E=S^1\times S^1$ i rzutu $E\to S^1$ oblicz stronę $E^2$ i $E^\infty$ homologicznego ciągu Serre'a. Wyjaśnij, dlaczego $d^2$ znikają z powodów wymiarowych.
\item Zastosuj ciąg Serre'a do wiązki $S^1\to S^3\to S^2$ i wskaż cztery niezerowe pozycje $E^2$. Dlaczego $d^2$ musi być izomorfizmem, jeśli znamy $H_1(S^3)$ i $H_2(S^3)$?
\item Uzasadnij, że układ współczynników $\mathcal H_1(S^1)$ dla wiązki opisującej butelkę Kleina jest nietrywialny. Oblicz jej $H_0,H_1,H_2$ z ciągu spektralnego.
\wskazowka{Dla współczynników ze znakiem $-$ na podstawie $S^1$ różniczka komórkowa ma postać $-2$.}
\item Dla $S^m\to S^m\times B\to B$ rozstrzygnij, czy generator $H^m(S^m)$ ma niezerową transgresję. Uzasadnij przy pomocy istnienia przekroju oraz klasy Eulera.
\item Rozpatrz zorientowaną wiązkę sferową $S^1\to E\to S^2$, której klasa Eulera jest równa $k$ razy generator $H^2(S^2;\Z)$. Dla $k\ne0$ wyznacz wszystkie grupy homologii całkowitej $E$.
\wskazowka{Na stronie $E^2$ różniczka $d^2:\Z\to\Z$ jest mnożeniem przez $k$. Uzyskaj $H_1(E)\cong\Z/|k|$, $H_2(E)=0$, $H_3(E)\cong\Z$.}
\item Udowodnij, że jeżeli podstawa jest CW-kompleksem o wymiarze $d$, to homologiczne różniczki Serre'a $d^r$ znikają dla $r>d$. Czy to wystarczy do wyznaczenia $H_*(E)$ bez problemu rozszerzeń?
\item Dla fibracji przestrzeni dróg $\Omega S^n\to PS^n\to S^n$, $n\ge2$, sprawdź, dlaczego pierwszą różniczką wychodzącą z $E^2_{n,0}$ mogącą być niezerową jest $d^n$.
\item Uzupełnij szczegóły dowodu tożsamości $(d^r)^2=0$ dla filtrowanego kompleksu i porównaj z dokładną parą.
\item Sformułuj precyzyjnie, kiedy lokalny system homologii włókna jest stały. Podaj przykład fibracji nad $S^1$, dla której $H_0(F)$ ma trywialną monodromię, ale $H_1(F)$ nie.
\end{enumerate}

\chapter{Eilenberg--Zilber i wzór Künnetha}
\begin{enumerate}
\item Oblicz $\mathrm{AW}$ na singularnym $1$-sympleksie $\sigma:\Delta^1\to X\times Y$ i sprawdź, że wzór uwzględnia oba końce.
\item Wyznacz znaki wszystkich $(2,1)$-przeplotów i oblicz $\mathrm{sh}(\sigma_2\otimes\tau_1)$.
\item Oblicz homologię tensoru kompleksów $P_6\otimes P_{10}$, gdzie $P_m=(\Z\xrightarrow{m}\Z)$ w stopniach $1,0$.
\wskazowka{Zastosuj $d=\gcd(6,10)$, ale sprawdź też osobno cykle i brzegi w stopniu $1$.}
\item Wykaż, że $H^3(\R P^2\times\R P^2;\Z)=\Z/2$ i wyjaśnij, dlaczego iloczyny kubkowe klas stopnia dodatniego nie generują tej grupy.
\item \textbf{Künneth bez skończonej generacji.}
Niech $A=\bigoplus_{j\ge1}\Z/2^j$. Oblicz
$A\otimes_\Z\Z/2$ oraz $\operatorname{Tor}^{\Z}_1(A,\Z/2)$
za pomocą wolnych rezolwent.
\wskazowka{Użyj sumy kompleksów
$\Z\xrightarrow{2^j}\Z$ i tensora z $\Z/2$.}
\end{enumerate}

\chapter{Homologia i kohomologia \v{C}echa}
\begin{enumerate}
\item Narysuj nerw trzech łuków na $S^1$ z pustym przecięciem potrójnym. Następnie zmień pokrycie tak, aby pojawił się wypełniony trójkąt.
\item Uzasadnij szczegółowo, że dwa wybory mapy nerwów przy rozdrobnieniu są przyległe, i zbuduj homotopię ich realizacji.
\item Oblicz granicę odwrotną układu $\Z^m\leftarrow\Z^{m+1}$, w którym mapa odrzuca ostatnią współrzędną. Porównaj sumę prostą i iloczyn przeliczalny.
\item Dla solenoidu $\Sigma_3$ oblicz $\check H_1(\Sigma_3;\Z)$ i $\check H^1(\Sigma_3;\Z)$.
\wskazowka{Pierwszy układ to $\Z\xleftarrow{\cdot3}\Z\xleftarrow{\cdot3}\cdots$, a drugi otrzymujesz po odwróceniu strzałek.}
\item Wyjaśnij, czemu $H_0^{\mathrm{sing}}(\Sigma_p;\Z)$ nie musi zgadzać się z $\check H_0(\Sigma_p;\Z)$.
\item Przypuść, że $A\subset S^3$ jest rozłączną sumą $r$ osadzonych okręgów. Oblicz $H_1(S^3\setminus A;\Z)$ za pomocą dualności Aleksandra.

\item \textbf{Realizacja nerwu.} Skonstruuj jawnie trzy otwarte łuki pokrywające $S^1$, dla których przecięcia podwójne są niepuste, a przecięcie potrójne puste. Oblicz pierwszą homologię realizacji nerwu.
\wskazowka{Wybierz łuki symetryczne o długości nieco większej niż $2\pi/3$, ale mniejszej niż $4\pi/3$; sprawdź, czy rzeczywiście przecięcie potrójne pozostaje puste.}
\item \textbf{Kierunek strzałek.} Dla rozdrobnienia pokryć $\mathcal V\succ\mathcal U$ wyprowadź kierunki map indukowanych w homologii i kohomologii nerwów. Wyjaśnij, dlaczego w definicjach \v{C}echa występują odmienne rodzaje granic.
\item \textbf{Granica okręgów.} Dla układu $S^1\xleftarrow{z\mapsto z^p}S^1\xleftarrow{z\mapsto z^p}\cdots$ policz granice grup $H_1$ i $H^1$ oraz uzasadnij otrzymanie $\Z[1/p]$ po stronie kohomologicznej.
\wskazowka{W homologii otrzymujesz układ $\Z\xleftarrow{\times p}\Z\leftarrow\cdots$, a w kohomologii jego odwrócenie.}

\end{enumerate}

\chapter{Lokalne współczynniki i dualności}
\begin{enumerate}
\item Dla systemu na $S^1$ o monodromii $-1$ na $\Z$ oblicz $H^0$, $H^1$, $H_0$, $H_1$; porównaj znaczenie niezmienników i coinwariantów.
\item Oblicz wszystkie grupy w dwóch wersjach dualności Poincarégo--Lefschetza dla cylindra $S^1\times[0,1]$.
\item Wykaż, że dopełnienie $r$ rozłącznych okręgów w $S^2$ ma $r+1$ składowych, korzystając wyłącznie z dualności Aleksandra.
\item Podaj przykład zwartego nieorientowalnego $M$ z brzegiem, dla którego $H_n(M,\partial M;\Z)=0$, ale $H_n(M,\partial M;\mathcal O_M)=\Z$.
\wskazowka{Wybierz wstęgę Möbiusa i porównaj monodromię zwykłych oraz skręconych współczynników.}

\item \textbf{Lokalny system na okręgu.} Niech $\mathcal L$ będzie lokalnym systemem z włóknem $\Z$ i monodromią $-1$ na $S^1$. Zapisz komórkowy kompleks z lokalnymi współczynnikami i oblicz $H_0(S^1;\mathcal L)$ oraz $H_1(S^1;\mathcal L)$.
\wskazowka{Różniczka w komórkowym kompleksie okręgu ma postać $T-\id$, gdzie $T=-1$.}

\end{enumerate}

\chapter{Przestrzenie klasyfikujące i teoria przeszkód}
\begin{enumerate}
\item \textbf{Przestrzenie $K(G,1)$.} Udowodnij, że iloraz
$S^\infty/(z\sim -z)$ jest $K(\Z/2,1)$. W jakich stopniach
jego grupy homotopii są niezerowe?
\wskazowka{Przestrzeń $S^\infty$ jest ściągalna, a rzutowanie jest nakryciem dwukrotnym.}
\item \textbf{Dwie konstrukcje $K(\Z,2)$.} Wykaż, że $\C P^\infty$
ma dokładnie jedną niezerową dodatnią grupę homotopii.
Porównaj argument z długim ciągiem fibracji $S^1\to S^\infty\to\C P^\infty$.
\item \textbf{Brak równoważności homotopijnej.}
Porównaj pierścienie kohomologii $\C P^2$ i $S^2\vee S^4$.
Dlaczego sam izomorfizm grup homologii nie wystarcza?
\wskazowka{Oblicz kwadrat generatora stopnia dwa.}
\item \textbf{Klasyfikacja nakryć.} Niech $X=S^1\vee S^1$ i $G=\Z/3$.
Opisz klasy izomorfizmu głównych wiązek $G$ nad $X$.
\wskazowka{Mamy $\pi_1(X)=F_2$; homomorfizmy $F_2\to\Z/3$
są wyznaczone przez dwa elementy grupy.}
\item \textbf{Monodromia i sprzężenie.}
Dla $G=S_3$ policz liczbę klas izomorfizmu głównych wiązek
$G$ nad $S^1$.
\wskazowka{Wiązki klasyfikują klasy sprzężoności elementów $S_3$.}
\item \textbf{Funkcje sklejające.} Rozpatrz $g_k:S^1\to U(1)$,
$g_k(e^{it})=e^{ikt}$. Porównaj wiązki otrzymane dla
$k=0,1,-1,2$, obliczając pierwszą klasę Cherna.
\item \textbf{Przekroje.} Dlaczego zespolona wiązka liniowa
$L_k\to S^2$ ma globalny nigdzie niezerowy przekrój
dokładnie wtedy, gdy $k=0$?
\wskazowka{Przekrój trywializuje wiązkę liniową; zastosuj klasyfikację przez $H^2(S^2;\Z)$.}
\item \textbf{Przeszkoda komórkowa.} Niech $X$ ma jedną
$0$-komórkę, jedną $n$-komórkę i jedną $(n+1)$-komórkę,
przyklejoną mapą stopnia $m$ na $S^n$, gdzie $n\ge2$.
Dla $A=\Z$ oblicz $H^n(X;A)$ oraz $H^{n+1}(X;A)$
i zinterpretuj obliczenie przez mapy do $K(\Z,n)$.
\wskazowka{Jedyna nietrywialna koróżniczka komórkowa ma postać
$\Z\xrightarrow{\,m\,}\Z$.}
\item \textbf{Sfera parzystowymiarowa.} Korzystając z klasy
Eulera i twierdzenia Poincarégo--Hopfa, udowodnij brak
nigdzie niezerowego pola stycznego na $S^{2m}$.
\item \textbf{Sfera nieparzystowymiarowa.} Skonstruuj
jawne nigdzie niezerowe pole styczne na $S^{2m+1}$.
\wskazowka{Zidentyfikuj $\R^{2m+2}$ z $\C^{m+1}$ i rozważ $x\mapsto ix$.}
\item \textbf{Zerowa klasa Eulera bez przekroju.}
Zbuduj nietrywialną orientowaną wiązkę rzędu trzy nad $S^4$
i udowodnij, że jej klasa Eulera jest zerowa, ale wiązka
nie ma nigdzie niezerowego przekroju.
\wskazowka{Weź niezerowy element $\pi_3(SO(3))\cong\Z$.
Przekrój rozkładałby wiązkę na liniową trywialną
oraz orientowaną dwuwymiarową; tymczasem $\pi_3(SO(2))=0$.}
\item \textbf{Pierwsza przeszkoda nie zawsze wystarcza.}
Sformułuj, w którym miejscu dowodu o przekroju
potrzebne są wyższe grupy homotopii włókna.
Podaj przykład włókna, które ma więcej niż jedną
niezerową grupę homotopii.
\wskazowka{Sfera $S^2$ ma zarówno $\pi_2\cong\Z$,
jak i $\pi_3\cong\Z$.}
\item \textbf{Ogólne fibracje Serre'a.}
Zapisz mapę $E\to E^h$ z dowodu ogólnego twierdzenia Serre'a i wykaż, że jest ona
równoważnością homotopijną przestrzeni całkowitych.
Uzasadnij, dlaczego ten argument nie daje sam przez się
równoważności na włóknach.
\wskazowka{Homotopia skracająca drogę nie zachowuje
rzutowania $p^h$ na podstawę.}
\end{enumerate}


\chapter{Wieże Postnikowa, iloczyny Whiteheada i operacje Steenroda}
\begin{enumerate}
\item \textbf{Etapy Postnikowa przestrzeni asferycznej.}
Niech $X=K(G,1)$ dla grupy dyskretnej $G$. Ustal wszystkie
etapy $P_nX$, $n\ge1$, w kategorii homotopijnej.
\wskazowka{Dla $j>1$ wszystkie grupy $\pi_j$ są już zerowe.}
\item \textbf{Etap drugiego rzędu.}
Udowodnij, że dla jednospójnego CW-kompleksu $X$
o $\pi_2(X)\cong A$ zachodzi $P_2X\simeq K(A,2)$.
Sprawdź ten wniosek dla $S^2$ i $\C P^2$.
\item \textbf{Włókno homotopijne.}
Dla mapy stałej $k:B\to K(A,n+1)$ oblicz typ homotopijny
$\operatorname{hofib}(k)$; wskaż, skąd bierze się
rozdzielenie na czynniki.
\wskazowka{Drogi kończące się w bazowym punkcie tworzą
model $\Omega K(A,n+1)$ nad mapą stałą.}
\item \textbf{Znikająca klasa przeszkody.}
Dla $Z=S^2\times S^2$ niech
$c=px+qy\in H^2(Z;\Z)$, gdzie $x,y$ są generatorami
pochodzącymi z rzutowań. Ustal, dla jakich $p,q\in\Z$
mapa $Z\to K(\Z,2)$ reprezentująca $c$ podnosi się
do $P_3S^2$.
\wskazowka{Oblicz $(px+qy)^2=2pq(x\smile y)$.
Uwaga: wynik dotyczy $P_3S^2$, a nie automatycznie
mapy do samej $S^2$.}
\item \textbf{Komórka dołączająca i Whitehead.}
Rozpisz CW-rozkład $S^p\times S^q$ dla $p,q\ge2$.
Uzasadnij, że mapa $S^p\vee S^q\to X$ przedłuża się
do iloczynu wtedy i tylko wtedy, gdy odpowiedni
iloczyn Whiteheada zanika.
\item \textbf{Niezmiennik Hopfa.}
Dla $f:S^3\to S^2$ i $k\in\Z$ porównaj
$H(kf)$ z $H(f)$, jeżeli $kf$ oznacza
$k$-krotność w $\pi_3(S^2)$. Wykaż, że
$H([\iota_2,\iota_2])=\pm2$.
\wskazowka{Porównaj kwadrat klasy $x+y$ na
$S^2\times S^2$ z kohomologią stożka mapy dołączającej.}
\item \textbf{Kołańcuch Steenroda.}
Dla $a\in Z^n(X;\F_2)$ i $0\le i\le n$,
wyprowadź z relacji dla $\smile_r$ równość
$\delta(a\smile_{n-i}a)=0$.
\item \textbf{Niezależność od reprezentanta.}
Jeśli $a'=a+\delta b$, wykaż bez skrótów, że
$(a'\smile_r a')-(a\smile_r a)$ jest koróżniczką.
\wskazowka{Rozwiń wyrażenie i zastosuj relację
$\delta(a\smile_r b)=\delta a\smile_r b+
 a\smile_r\delta b+a\smile_{r-1}b+b\smile_{r-1}a$.}
\item \textbf{Kwadraty Steenroda na przestrzeni rzutowej.}
Niech $t\in H^1(\R P^\infty;\F_2)$.
Oblicz $\operatorname{Sq}^i(t^5)$ dla $0\le i\le5$.
Następnie powtórz rachunek na $\R P^6$.
\wskazowka{Współczynniki są $\binom5i$ modulo $2$;
uwzględnij ograniczenie wymiaru przestrzeni.}
\item \textbf{Działanie na uniwersalnej klasie Cherna.}
Niech $u\in H^2(\C P^\infty;\F_2)$.
Oblicz $\operatorname{Sq}^2(u^3)$ i
$\operatorname{Sq}^4(u^3)$.
\wskazowka{Skorzystaj ze wzoru Cartana i
$\operatorname{Sq}^2(u)=u^2$.}
\item \textbf{Relacja Adema w najniższym wymiarze.}
Zastosuj podany wzór z $a=b=1$ i oblicz
$\operatorname{Sq}^1\operatorname{Sq}^1$.
Dlaczego sama naturalność i stopień operacji nie
wystarczają do wyprowadzenia wszystkich relacji Adema?
\item \textbf{Lokalne współczynniki w wieży.}
Wyjaśnij, dlaczego w ogólnej przestrzeni z
niezerową $\pi_1(X)$ działania grupy fundamentalnej
na $\pi_n(X)$ nie można pominąć w definicji
niezmienników $k$. Podaj przestrzeń,
której grupa fundamentalna działa niebanalnie
na wyższej grupie homotopii.
\wskazowka{Rozważ nietrywialną wiązkę $S^2$
nad $S^1$ z monodromią stopnia $-1$.}
\item \textbf{Zanik pierwszej przeszkody nie jest końcem.}
Porównaj możliwość podniesienia mapy
$Z\to P_2S^2$ do $P_3S^2$ i dalej do
$P_4S^2$. Jakie dane poza samą grupą
$\pi_3(S^2)$ musimy znać, aby kontynuować?
\end{enumerate}

\chapter{Operacje kohomologiczne i relacje Adema}
\begin{enumerate}
\item \textbf{Wykonawcza formuła Bocksteina.}
Niech $X=\R P^2$. Wykaż, że $\beta_2(t)=t^2$ dla
$t\in H^1(X;\F_2)$. Sformułuj wynik przy użyciu
konkretnej podniesionej koróżniczki komórkowej.
\wskazowka{Komórkowa różniczka w stopniu dwa to mnożenie przez $2$;
współczynniki modulo $4$ nie usuwają tej różniczki.}
\item \textbf{Działanie na produkcie rzutowym.}
Oblicz $\operatorname{Sq}^2(u^3v^2)$ w
$H^*(\R P^6\times\R P^3;\F_2)$ i wskaż, które
składniki znikają z powodów wymiarowych.
\wskazowka{Rozpisz przypadki $(i,j)=(0,2),(1,1),(2,0)$
w formule Cartana.}
\item \textbf{Relacje Adema na konkretnym jednomianie.}
Porównaj $\operatorname{Sq}^2\operatorname{Sq}^2(t^m)$
i $\operatorname{Sq}^3\operatorname{Sq}^1(t^m)$ w
$H^*(\R P^\infty;\F_2)$ dla $m=1,2,3,4,5$.
Czy takie sprawdzenie dowodzi relacji na dowolnej przestrzeni?
\wskazowka{Użyj $\operatorname{Sq}^i(t^m)=\binom mi t^{m+i}$
i odróżnij przykład od tożsamości operacji naturalnych.}
\item \textbf{Słowa dopuszczalne.}
Korzystając z relacji Adema, zapisz
$\operatorname{Sq}^1\operatorname{Sq}^2$ oraz
$\operatorname{Sq}^2\operatorname{Sq}^3$ jako sumy złożeń
dopuszczalnych. Wyjaśnij, dlaczego pierwszy iloczyn
jest jedną operacją, a nie sumą wielu wyrazów.
\item \textbf{Kwadraty klasy Thoma.}
Niech $\xi\to B$ będzie rzeczywistą wiązką rzędu $r$
ze znormalizowaną klasą Thoma $U$ nad $\F_2$.
Wykaż, że $w_i(\xi)$ jest jednoznacznie wyznaczona
przez $\operatorname{Sq}^i(U)$ i że $w_i=0$ dla $i>r$.
\wskazowka{Wykorzystaj izomorfizm Thoma i niestabilność operacji.}
\item \textbf{Rozszczepienie wiązki tautologicznej.}
Wyznacz wszystkie klasy $w_i(\gamma^\perp)$ nad
$\R P^4$ i sprawdź, że całkowita klasa wiązki
$\gamma\oplus\gamma^\perp$ jest równa $1$.
\item \textbf{Pierścień bez iloczynów i nietrywialny kwadrat.}
Oblicz grupy kohomologii i wszystkie iloczyny dodatnich
stopni dla $\Sigma\C P^2$. Następnie oblicz
$\operatorname{Sq}^2:H^3(\Sigma\C P^2;\F_2)\to
H^5(\Sigma\C P^2;\F_2)$.
\wskazowka{Użyj zredukowanego izomorfizmu zawieszenia i
$\operatorname{Sq}^2(u)=u^2$ na $\C P^2$.}
\item \textbf{Sama naturalność nie wystarcza.}
Znajdź naturalne odwzorowanie $H^2(X;\Z)\to H^4(X;\Z)$,
które w ogólności nie jest addytywne. Wskaż przestrzeń,
na której addytywność zawodzi.
\wskazowka{Rozważ $u\mapsto u^2$ na $\C P^\infty\times\C P^\infty$.}
\end{enumerate}

\chapter{Zawieszenie i homotopia stabilna}
\begin{enumerate}
\item \textbf{Stopień i zawieszenie.}
Niech $f:S^n\to S^n$ ma stopień $d$, $n\ge1$.
Udowodnij, że $\deg(\Sigma f)=d$.
\wskazowka{Porównaj mapę na zredukowanych homologiach
z izomorfizmem zawieszenia.}
\item \textbf{Zakres stabilizacji.}
Ustal najmniejsze $n$ gwarantowane przez twierdzenie
Freudenthala, dla którego
$\pi_{n+4}(S^n)\to\pi_{n+5}(S^{n+1})$
jest izomorfizmem. Czy twierdzenie zapewnia minimalność
tej wartości wśród wszystkich możliwych $n$?
\wskazowka{Dla $S^n$ izomorfizm zachodzi, gdy $n+4<2n-1$.}
\item \textbf{Ostry zakres dla $S^2$.}
Korzystając z równania $[\iota_2,\iota_2]=\pm2\eta$,
wykaż, że $E(2\eta)=0$. Uzasadnij,
dlaczego nie wynika stąd jeszcze $E\eta\ne0$.
\wskazowka{Zawieszenie iloczynu Whiteheada zanika.}
\item \textbf{Pierwszy stabilny pień sfer.}
Odtwórz dowód $\pi_1^{\mathrm s}=\Z/2$,
pisząc osobno: surjektywność Freudenthala,
relację Whiteheada, niezerowość wykrytą przez $\operatorname{Sq}^2$.
\item \textbf{Stożek zawieszonej mapy.}
Dla $f:S^p\to S^q$, $p>q\ge1$, skonstruuj
homotopijną równoważność $C_{\Sigma f}\simeq\Sigma C_f$.
\wskazowka{Zapisz oba stożki jako odpowiednie ilorazy
$S^q\sqcup(S^p\times[0,1])$ i porównaj relacje.}
\item \textbf{Zero i jedynka widma sferycznego.}
Wyprowadź $\pi_0^{\mathrm s}=\Z$ i opisz
reprezentanta generatora. Wyjaśnij, dlaczego dla
$k=1$ nie można pozostać przy grupie $\pi_3(S^2)$.
\item \textbf{Widmo Eilenberga--Mac Lane'a.}
Dla $A=\Z/6\Z$ oblicz $\pi_k(HA)$ dla wszystkich
$k\in\Z$. Wyjaśnij, dlaczego obliczenie nie mówi,
że wszystkie przestrzenie $K(A,n)$ są ściągalne poza
jednym wybranym wymiarem.
\item \textbf{Dlaczego potrzebujemy widm?}
Wybierz $n\ge2$ i porównaj
$\pi_{n+1}(S^n)$ oraz $\pi_{n+1}(K(\Z,n))$.
Jaka informacja o stabilnych grupach ginie po
zastąpieniu widma sferycznego widmem $H\Z$?
\wskazowka{Dla $n\ge3$ skorzystaj z pierwszego
stabilnego pnia oraz definicji $K(\Z,n)$.}
\item \textbf{Granice twierdzeń.}
Sformułuj, jakie dodatkowe twierdzenie jest użyte w dowodzie
Freudenthala, a jakie w dowodzie niezerowości $E\eta$.
Oddziel wyniki udowodnione w skrypcie od twierdzeń
przyjętych jako narzędzia pomocnicze.
\end{enumerate}

\chapter{Pontriagin--Thom i teoria kobordyzmów}
\begin{enumerate}
\item \textbf{Wartość regularna i normalne uramowanie.}
Niech $f:S^3\to S^2$ będzie gładką mapą, a $q$ jej
wartością regularną. Wykaż, że $f^{-1}(q)$ jest skończoną
sumą okręgów i opisz otrzymane uramowanie normalne.
\wskazowka{Skorzystaj z twierdzenia o wartości regularnej,
klasyfikacji zwartych $1$-rozmaitości i ilorazu
$T_xS^3/\ker\mathrm{D} f_x\cong T_qS^2$.}
\item \textbf{Dlaczego trzeba zapamiętać uramowanie?}
Podaj dwa uramowania wiązki normalnej standardowego
okręgu w $\R^3$ różniące się o jeden pełny obrót.
Wyjaśnij, jak różnica ta może być niewidoczna po
zapomnieniu uramowania.
\wskazowka{Porównaj klasy map $S^1\to SO(2)$ i odróżnij
izomorfizm wiązki od konkretnej trywializacji.}
\item \textbf{Trywialna wiązka i przestrzeń Thoma.}
Udowodnij $\operatorname{Th}(B\times\R^k)\cong B_+\wedge S^k$.
Oblicz $\widetilde H^*(\operatorname{Th}(B\times\R^k);\Z)$
dla $B=S^2$, przy $k\ge1$.
\wskazowka{Skorzystaj z definicji ilorazu
$B\times D^k/B\times S^{k-1}$ i izomorfizmu zawieszenia.}
\item \textbf{Ciągłość odwzorowania ściągającego.}
Napisz jawny wzór mapy $S^m\to S^k$ otrzymywanej
z trywialnego otoczenia tubularnego uramowanej
podrozmaitości $M^{m-k}\subset S^m$, i wyjaśnij
ciągłość na brzegu otoczenia.
\wskazowka{Na dysku normalnym użyj ilorazu
$D^k/S^{k-1}$; dopełnienie otoczenia wyślij w punkt bazowy.}
\item \textbf{Stopień i punkty uramowane.}
Dla $f(z)=z^d$ na $S^1\subset\C$, $d\in\Z$,
oblicz liczbę punktów $f^{-1}(1)$ ze znakami
pochodnej i porównaj ją ze stopniem $f$.
Omów przypadek $d=0$ osobno.
\item \textbf{Uramowanie z pochodnej.}
Wyjaśnij szczegółowo, dlaczego wybrana baza
$T_qS^k$ daje bazę normalną do $f^{-1}(q)$,
ale nie daje w ogólności trywializacji $T(f^{-1}(q))$.
\item \textbf{Uramowany okrąg Hopfa.}
Wyjaśnij, dlaczego regularne włókno odwzorowania
Hopfa daje niezerową klasę w $\Omega_1^{\mathrm{fr}}$,
ale jest zerem w $\Omega_1^O$.
\wskazowka{Porównaj wynik $\pi_1^{\mathrm s}\cong\Z/2$
z faktem, że okrąg jest brzegiem dysku.}
\item \textbf{Grupa kobordyzmu w wymiarze zero.}
Wyznacz $\Omega_0^O$ i $\Omega_0^{SO}$,
konstruując jawne kobordyzmy, które identyfikują
pary punktów (bez orientacji oraz z orientacjami
przeciwnymi).
\item \textbf{Produkt i orientacja.}
Niech $M^m,N^n$ będą zamkniętymi zorientowanymi
rozmaitościami. Sprawdź w lokalnych współrzędnych,
że zamiana czynników na $M\times N$ zmienia
orientację o $(-1)^{mn}$.
\item \textbf{Zanik liczb charakterystycznych na brzegu.}
Odtwórz dowód twierdzenia o zaniku liczb Stiefela--Whitneya na brzegu bez
odwołań do jego treści: rozpisz stabilny rozkład
$TW|_{\partial W}$, długi ciąg pary $(W,\partial W)$
i naturalność ewaluacji klas kohomologii.
\wskazowka{Najpierw udowodnij
$i_*[\partial W]_2=0$, a potem przejdź do klas $w_j$.}
\item \textbf{Klasy charakterystyczne $\R P^2$.}
Z $T\R P^2\oplus\underline\R\cong3\gamma^1$
wyprowadź $w_1$ i $w_2$; oblicz oba stopnia $2$
i wyjaśnij, dlaczego $\R P^2$ nie może być brzegiem.
\item \textbf{Kobordyzm sumy spójnej.}
Naszkicuj trójwymiarowy kobordyzm od
$M\sqcup N$ do $M\#N$ dla dwóch zamkniętych
powierzchni. Wyjaśnij, w którym miejscu dołączamy
jednowymiarowy uchwyt.
\item \textbf{Pierwsza nieorientowana grupa dwuwymiarowa.}
Przyjmując klasyfikację zamkniętych powierzchni,
udowodnij $\Omega_2^O\cong\Z/2$ i sprawdź,
że izomorfizm jest zadany przez $\chi\bmod2$.
\wskazowka{Sumy spójne zamień na sumy rozłączne
w kobordyzmie; niezerowość wykryj przez $w_1^2$.}
\item \textbf{Uramowanie a klasyfikacja wiązek.}
Niech $\nu\to M$ będzie wiązką normalną rzędu $k$.
Wykaż, że uramowanie $\nu$ daje stałą mapę
klasyfikującą $M\to BO(k)$, ale odwrotna implikacja
wymaga uwzględnienia konkretnej trywializacji.
\item \textbf{Od $MO(k)$ do rozmaitości.}
Niech $f:S^{n+k}\to\operatorname{Th}(\gamma_{k,N})$
będzie mapą poprzeczną do przekroju zerowego.
Oblicz wymiar jego przeciwobrazu i określ
wiązkę normalną z wykorzystaniem $f$.
\wskazowka{Sekcja zerowa ma kowymiar $k$ w całej
przestrzeni wiązki; użyj twierdzenia o przeciwobrazie
transwersalnym.}
\item \textbf{Dualność przez ściągnięcie.}
Niech $M\subset Q$ będzie zamkniętą, zorientowaną
podrozmaitością kowymiaru $k$ w zamkniętej
zorientowanej rozmaitości. Sprawdź w lokalnym modelu
$\R^{m-k}\times\R^k$ równość
$\operatorname{PD}_Q[M]=c_{M,Q}^*U$.
\wskazowka{Użyj lokalnego generatora
$H^k(D^k,S^{k-1};\Z)$ i znaku przecięcia.}
\item \textbf{Zera przekrojów.}
Wykaż, że dla dowolnej zorientowanej rzeczywistej
wiązki liniowej nad $S^1$ klasa Eulera znika.
Porównaj z pierwszą klasą Stiefela--Whitneya
wiązki Möbiusa, która nie jest zorientowana.
\item \textbf{Co jest twierdzeniem, a co redukcją?}
Rozdziel dowody zawarte w skrypcie od przyjętych
narzędzi pomocniczych w konstrukcji Pontriagina--Thoma.
Wymień dokładne miejsca wykorzystania twierdzeń
Whitneya, transwersalności względnej, tubularności
i skończenie wymiarowej aproksymacji $BO(k)$.
\end{enumerate}

\chapter{Homologia Morse'a: trajektorie i komórki}
\begin{enumerate}
\item \textbf{Hesjan w punkcie krytycznym.}
Pokaż, że macierz drugich pochodnych funkcji gładkiej w punkcie
krytycznym opisuje formę dwuliniową niezależną od wyboru mapy.
\wskazowka{Zastosuj regułę łańcuchową drugiego rzędu; wszystkie dodatkowe
składniki zawierają pierwszą pochodną funkcji.}
\item \textbf{Indeksy na torusie.} Dla $f(\theta,\varphi)=\cos\theta+2\cos\varphi$
na $T^2$ wyznacz punkty krytyczne, hesjany i indeksy. Uzasadnij,
że funkcja jest funkcją Morse'a.
\item \textbf{Punkty krytyczne sfery.} Na $S^n$ rozważ wysokość
$x\mapsto x_{n+1}$. Zdefiniuj indeksy i napisz kompleks
Morse'a dla $n=1,2,3$, wyjaśniając przypadek $n=1$.
\item \textbf{Rachunek energii.} Niech $\gamma:\R\to M$ łączy
punkty krytyczne $p,q$ ujemnym przepływem gradientowym.
Udowodnij $\int_{t\in\R}\|\operatorname{grad}f(\gamma(t))\|^2_g\,\dd t=f(p)-f(q)$.
Wyjaśnij, dlaczego nie ma nietrywialnych trajektorii od $p$ do $q$
przy $f(p)\le f(q)$.
\item \textbf{Wymiar przestrzeni trajektorii.} Zapisz
wymiary $W^u(p)$, $W^s(q)$ i ich przekroju. Uzasadnij,
że przekrój niesparametryzowanych trajektorii z regularnym
poziomem $f=c$ obniża wymiar o jeden.
\item \textbf{Dlaczego $\partial^2=0$?}
Dla punktów krytycznych o indeksach $k,k-1,k-2$ odtwórz
pełny argument z brzegiem jednowymiarowej zwartej
kompaktyfikacji $\overline{\mathcal M}(p,r)$.
Wyraźnie wskaż, gdzie używa się zewnętrznego twierdzenia
o sklejaniu trajektorii.
\wskazowka{Zinterpretuj każdy iloczyn $n(p,q)n(q,r)$ jako
sumę znaków punktów brzegowych; odcinki mają po dwa końce
z przeciwnymi znakami.}
\item \textbf{Niedoskonały okrąg.}
Dla $f(\theta)=\cos(2\theta)$ na $S^1$ wybierz
orientacje niestabilnych kierunków i oblicz macierz
różniczki Morse'a. Wskaż jawne generatory homologii.
\item \textbf{Zmiana orientacji.} Zmień orientację $W^u(p)$
w jednym punkcie krytycznym. Pokaż, że nowy kompleks
jest izomorficzny ze starym przez pomnożenie generatora $p$
przez $-1$; uwzględnij różniczki przychodzące i wychodzące.
\item \textbf{Podpoziomy.} Rozważ funkcję Morse'a
$f:S^2\to\R$ będącą wysokością. Opisz typy homotopijne
jej podpoziomów poniżej minimum, między punktami krytycznymi
i powyżej maksimum.
\item \textbf{Nierówności Morse'a.} Wyprowadź tożsamość
$m_k=b_k+r_k+r_{k+1}$ i silne nierówności Morse'a.
Wykaż, że minimalna możliwa liczba punktów krytycznych
na torusie wynosi cztery.
\item \textbf{Torsja a liczba punktów krytycznych.} Zastanów się,
co można powiedzieć o minimalnej liczbie generatorów kompleksu Morse'a
w obecności niezerowej torsji w homologii całkowitej.
\wskazowka{Rozważ krótki wolny kompleks $\Z\xrightarrow{m}\Z$.
Rangi Betti nie wykrywają jego torsji, lecz liczby komórek już tak.}
\item \textbf{Jednakowe wartości krytyczne.}
Niech $f(\theta,\varphi)=\cos\theta+\cos\varphi$ na torusie
$T^2$. Pokaż, że dwa punkty siodłowe mają tę samą wartość
krytyczną. Wyznacz względną homologię
$H_j(M^{\varepsilon},M^{-\varepsilon};\Z)$ dla małego
$\varepsilon>0$, gdzie $M^a=f^{-1}(( -\infty,a])$.
\wskazowka{Przy przejściu przez poziom zero dołączamy
jednocześnie dwa uchwyty indeksu $1$; użyj ekscyzji
i względnej homologii sumy dwóch par $(D^1,S^0)$.}
\item \textbf{Niezależność znaków od orientowalności $M$.}
Uzasadnij, dlaczego do określenia znaków sztywnych trajektorii
Morse'a wystarczy orientować niestabilne rozmaitości $W^u(p)$,
a nie całą rozmaitość $M$.
\wskazowka{Orientacja $W^u(q)$ orientuje normalną do $W^s(q)$;
porównaj ją z orientacją $W^u(p)$ i kierunkiem przepływu.}
\item \textbf{Kontynuacja.} Wypisz stopnie i warunki końcowe
dla trajektorii parametryzowanej rodziny $(f_s,g_s)$, która
ma definiować mapę łańcuchową. Wskaż, jaka
kompaktyfikacja jest potrzebna do dowodu zgodności z różniczką.
\end{enumerate}

\chapter{Indeks i twierdzenie Lefschetza}
\begin{enumerate}
\item \textbf{Ślad Hopfa.} Dla kompleksu
$0\to\Q^2\xrightarrow{\mathrm{diag}(1,0)}\Q^2\to0$
wybierz endomorfizm łańcuchowy dany dwiema diagonalnymi
macierzami. Oblicz naprzemienny ślad na łańcuchach
i na homologii, weryfikując Lemat śladu Hopfa.
\item \textbf{Homotopie.} Udowodnij, że $L(f)$ jest
niezmiennikiem klasy homotopii odwzorowania $f:X\to X$
dla skończonego kompleksu CW $X$.
\item \textbf{Twierdzenie Brouwera z Lefschetza.}
Podaj dowód istnienia punktu stałego ciągłej mapy
$D^n\to D^n$, podając wyraźnie wszystkie hipotezy.
\item \textbf{Dlaczego trzeba dzielić sympleksy?}
Załóż, że $f:|K|\to|K|$ nie ma punktów stałych.
Odtwórz argument wykorzystujący $\delta=\min_x\|f(x)-x\|>0$,
subdywizję i aproksymację symplicjalną.
Wykaż, że w odpowiedniej mapie łańcuchowej
wszystkie elementy diagonalne mają współczynnik zero.
\wskazowka{Porównaj $f(\sigma)$ z $\sigma$ dla sympleksu
o średnicy mniejszej niż $\delta/4$.}
\item \textbf{Przejście od wielościanów do CW.}
Załóż, że $X$ jest zwartym retraktem otoczenia $U$ w
$\R^N$. Skonstruuj wielościan $P\subset U$ zawierający $X$
i mapę $F:P\to P$ tak, by brak punktów stałych $f:X\to X$
pociągał brak punktów stałych $F$. Wykorzystaj
$\operatorname{tr}(AB)=\operatorname{tr}(BA)$.
\item \textbf{Sfery.} Oblicz $L(f)$ dla dowolnej mapy
$f:S^n\to S^n$ stopnia $d$. Dla których $d$ niezerowość
Lefschetza wymusza punkt stały? Porównaj z mapą antypodalną.
\item \textbf{Sfera stopnia minus jeden.}
Skonstruuj mapę $S^2\to S^2$ stopnia $-1$,
która ma punkt stały, oraz mapę tego samego stopnia,
która nie ma punktów stałych.
\wskazowka{Odbicie względem płaszczyzny ma punktów stałych wiele,
a antypodalna mapa nie ma żadnego.}
\item \textbf{Torus i wyznacznik.} Dla macierzy
$A=\begin{psmallmatrix}2&1\\1&1\end{psmallmatrix}$
wyznacz liczbę Lefschetza mapy $T^2\to T^2$ oraz
jej punkty stałe i ich lokalne indeksy.
\item \textbf{Indeks lokalny.} Niech $F:\R^2\to\R^2$
będzie liniowym odwzorowaniem z macierzą $A$
i załóż, że $\det(I-A)\ne0$. Oblicz indeks zera mapy
$\id-F$ z definicji stopnia. Sprawdź wynik dla
$A=\operatorname{diag}(2,1/2)$.
\item \textbf{Diagonala.} Dla zamkniętej zorientowanej
rozmaitości $M$ sprawdź, że warunek
$\det(I-\mathrm Df_p)\ne0$ jest równoważny
transwersalności $\Gamma_f$ i $\Delta$ w $(p,p)$.
\item \textbf{Znak diagonalnej.}
Dla $M=S^1$ i generatora $u\in H^1(S^1;\Z)$
ustal znak klasy $\operatorname{PD}_{T^2}[\Delta]$
w $H^1(S^1\times S^1;\Z)$, tak by lokalny indeks
mapy stałej był $+1$.
\item \textbf{Działanie torusa w homologii.}
Uzasadnij tożsamość
$\sum_{k=0}^d(-1)^k\operatorname{tr}(\bigwedge^kA)
=\det(I-A)$; możesz skorzystać z wielomianu charakterystycznego.
\item \textbf{Lefschetz identyczności.} Dla powierzchni
o rodzaju $g$ oblicz $L(\id)$ i porównaj z charakterystyką
Eulera otrzymaną ze struktury CW. Dlaczego dla torusa
wartość ta nie daje informacji o istnieniu punktów stałych?
\item \textbf{Zmiana współrzędnych w indeksie lokalnym.}
Niech $F:U\to\R^n$ ma izolowany punkt stały w $0$, a
$h:U\to h(U)$ będzie dyfeomorfizmem klasy $C^1$ z $h(0)=0$.
Uzasadnij równość indeksów $\operatorname{ind}_0(F)=
\operatorname{ind}_0(h\circ F\circ h^{-1})$, nawet gdy $F$ jest
tylko ciągłe.
\wskazowka{Zastosuj wzór $h(x)-h(F(x))=A(x)(x-F(x))$
ze $A(x)\to\D h_0$; uwzględnij znaki obu zmian stopnia.}
\item \textbf{Rozmaitość niespójna.}
Niech $M=S^2\sqcup S^2$ i niech $f:M\to M$ zamienia obie
składowe za pomocą ustalonych homeomorfizmów. Wyznacz $L(f)$,
obliczając bloki diagonalne na homologii, i sprawdź, że mapa
nie ma punktów stałych.
\item \textbf{Lokalne współczynniki.} Wyjaśnij, jakich
modyfikacji wymaga dowód wzoru indeksowego na rozmaitości
nieorientowalnej. Nie zakładaj istnienia globalnej klasy
fundamentalnej w homologii całkowitej bez skręcenia.
\wskazowka{Wróć do systemu orientacji z rozdziału o dualności.}
\end{enumerate}

\chapter{Analityczne podstawy homologii Morse'a}
\begin{enumerate}
\item \textbf{Punkty graniczne.}
Dla gładkiego przepływu gradientowego na zwartej rozmaitości
wykaż, że zbiór graniczny w przód jest spójny i zawiera się
w zbiorze krytycznym. Dlaczego teza nie musi zachodzić dla
przepływu bez funkcji Lapunowa?
\wskazowka{Użyj przecięcia zagnieżdżonych zwartych domknięć
ogonów trajektorii oraz monotoniczności funkcji Morse'a.}
\item \textbf{Całkowity czas z dala od punktów krytycznych.}
Niech $X=-\nabla f$, $f$ Morse'a na zwartej $M$.
Dla małych kul $U_p$ wokół punktów krytycznych
udowodnij jednostajne oszacowanie miary czasów, dla których
trajektoria leży w $M\setminus\bigcup_pU_p$.
\wskazowka{Zastosuj $\|\nabla f\|\ge m>0$ i bilans energii.}
\item \textbf{Granica łamana.}
Uzasadnij, dlaczego dla pary Morse'a--Smale'a
trajektoria łamana od $p$ do $q$ może mieć co najwyżej
$\lambda(p)-\lambda(q)$ nietrywialnych części.
\item \textbf{Sztywne trajektorie.}
Wyprowadź skończoność $\mathcal M(p,q)$, gdy
$\lambda(p)-\lambda(q)=1$, na podstawie zwartości
z dokładnością do łamania. Wskaż miejsce, w którym
potrzebna jest hipoteza Morse'a--Smale'a.
\item \textbf{Równanie Lapunowa.}
Jeżeli wszystkie wartości własne rzeczywistej macierzy $A$
mają ujemne części rzeczywiste, sprawdź, że
$P=\int_{t\in[0,\infty)}e^{tA^\top}e^{tA}\,\mathrm dt$
jest dodatnio określona i $A^\top P+PA=-I$.
\wskazowka{Zróżniczkuj iloczyn macierzy pod całką i
wykorzystaj zanik wykładniczy.}
\item \textbf{Gdzie potrzebne jest sklejanie?}
Wyjaśnij, dlaczego sam bilans energii i zbieżność ODE na
zwartych przedziałach nie dowodzą, że trajektoria łamana
jest końcem jednowymiarowej przestrzeni moduli.
Jaką rolę pełni jednostajna prawa odwrotność operatora
$\tfrac{\mathrm d}{\mathrm dt}-\mathrm DX(\gamma)$?

\item \textbf{Warunkowy argument sklejania.} Załóż, że $\|\mathcal F(\gamma_R)\|\le Ce^{-cR}$ i istnieją prawe odwrotne $Q_R$ operatorów linearyzacji, jednostajnie ograniczone. Wypisz stałe pozwalające zastosować zasadę kontrakcji do równania korekcyjnego.
\wskazowka{Oddziel oszacowanie liniowe od kwadratowej reszty Taylora; nie próbuj dowodzić samego istnienia $Q_R$ z oszacowania residuum.}

\end{enumerate}

\chapter{Indeks pól wektorowych i brzeg rozmaitości}
\begin{enumerate}
\item \textbf{Indeks pola w wymiarze dwa.}
Dla pól $X(x,y)=(x,y)$, $Y(x,y)=(x,-y)$ i
$Z(x,y)=(x^2-y^2,2xy)$ na $\R^2$ oblicz indeks zera
w początku. Uzasadnij obliczenie poprzez stopień mapy na $S^1$.
\wskazowka{Pierwsze dwa pola mają izolowane zera nieosobliwe.
Trzecie pole jest zapisem zespolonego odwzorowania $z\mapsto z^2$.}
\item \textbf{Zera gradientu i charakterystyka Eulera.}
Dla funkcji wysokości na $S^n$ oblicz indeks pola
gradientowego w minimum i maksimum. Sprawdź, że suma wynosi
$1+(-1)^n$.
\wskazowka{Hesjan w maksimum ma $n$ ujemnych wartości własnych.}
\item \textbf{Brzeg przedziału.}
Na $M=[0,1]$ skonstruuj pole skierowane ściśle na zewnątrz
na obu końcach. Wykaż, że suma indeksów zer wewnątrz wynosi $1$.
Czy pole stałe spełnia hipotezy twierdzenia?
\item \textbf{Podwojenie.}
Sprawdź, że podwojenie $D^n$ jest homeomorficzne z $S^n$.
Zdefiniuj składanie $r:S^n\to D^n$ dwóch półkul i wykaż,
że $r\circ j=\id_{D^n}$ dla naturalnego osadzenia $j$.
\item \textbf{Przejście od pola do mapy.}
Dla nieosobliwego zera $p$ pola $X$ i małego $t>0$
uzasadnij w lokalnych współrzędnych tożsamość
$\D_p\exp_x(-tX(x))=I-t\D_pX$.
\wskazowka{Różniczka mapy $(x,v)\mapsto\exp_x(v)$ w $(p,0)$
jest sumą składowych poziomej i pionowej.}
\end{enumerate}

\chapter{Kohomologia z nośnikami zwartymi i Borel--Moore'a}
\begin{enumerate}
\item \textbf{Definicja funktorialna.}
Dla właściwej mapy $f:X\to Y$ wypisz jawnie mapę par
$(X,X\setminus f^{-1}(K))\to(Y,Y\setminus K)$ i wyjaśnij,
dlaczego otrzymujemy $f^*:H_c^k(Y;G)\to H_c^k(X;G)$.
Znajdź niewłaściwą mapę, dla której ten argument nie działa.
\item \textbf{Pierwotna o nośniku zwartym.}
Załóż, że $f\in C_c^\infty(\R)$.
Udowodnij, że istnieje $F\in C_c^\infty(\R)$ spełniająca
$F'=f$ wtedy i tylko wtedy, gdy
$\int_{x\in\R} f(x)\,\mathrm dx=0$.
\item \textbf{Łańcuch prostej.}
Niech $\sigma_m(t)=m+t$, $m\in\Z$.
Wykaż, że $\sum_{m\in\Z}\sigma_m$ jest lokalnie
skończonym cyklem. Dlaczego ten sam wzór nie należy
do zwykłego kompleksu singularnego?
\item \textbf{Dwie grupy na zbiorze dyskretnym.}
Dla $X=\N$ z topologią dyskretną oblicz $H_0(X;\Z)$,
$H_0^{\mathrm{BM}}(X;\Z)$ i $H_c^0(X;\Z)$.
Porównaj z $H_0(X^+,\{\infty\};\Z)$.
\wskazowka{Na zbiorze dyskretnym każda zwarta podprzestrzeń jest skończona.}
\item \textbf{Współhomologia $\R^n$.}
Przeprowadź indukcję dla $H_{c,\mathrm{dR}}^*(\R^n)$,
sprawdzając na formach $\alpha+\mathrm dt\wedge\beta$ tożsamość
$\mathrm dH+H\mathrm d=\id-sI$.
\item \textbf{Otwarty cylinder.}
Oblicz $H_c^k(S^1\times\R;\Z)$ oraz
$H_k^{\mathrm{BM}}(S^1\times\R;\Z)$.
Podaj gładkie formy o nośniku zwartym reprezentujące
niezerowe klasy w stopniach $1$ i $2$ nad $\R$.
\item \textbf{Skończony i nieskończony nośnik.}
Niech $K$ będzie lokalnie skończoną triangulacją
niezw artej rozmaitości. Wyjaśnij, dlaczego komórkowe
kołańcuchy o nośniku zwartym mają skończenie wiele
niezerowych współczynników, natomiast zwykłe kołańcuchy
mogą mieć nieskończenie wiele współczynników.
\item \textbf{Wersja dualności z brzegiem.}
Dla $W=D^n$ i $M=\operatorname{int}W$ oblicz
$H^k(W,\partial W;\Z)$ oraz $H_k(W,\partial W;\Z)$.
Porównaj z grupami $H_c^k(M;\Z)$ i $H_k^{\mathrm{BM}}(M;\Z)$.
\item \textbf{Orientacje lokalne.}
Dla wstęgi Möbiusa bez brzegu (wnętrza zwartej wstęgi)
wskaż, gdzie w dualności Poincarégo konieczny jest
lokalny system orientacji.
\wskazowka{Transport orientacji wokół pętli centralnej zmienia znak.}
\item \textbf{Właściwość w geometrii algebraicznej.}
Wyjaśnij, dlaczego zanurzenie domknięte $i:Z\hookrightarrow X$
jest właściwe dla przestrzeni lokalnie zwartych Hausdorffa,
a otwarte zanurzenie na ogół nie jest. Sformułuj wnioski
dla $i_*:H_k^{\mathrm{BM}}(Z)\to H_k^{\mathrm{BM}}(X)$.
\end{enumerate}

\chapter{Kohomologia snopów z nośnikami i dualność Verdiera}
\begin{enumerate}
\item \textbf{Snop i przedsnop.}
Niech $X=\{a,b\}$ z topologią dyskretną. Wykaż, że przedsnop
przyporządkowujący każdemu niepustemu otwartemu podzbiorowi $k$
i każdemu niepustemu zawieraniu identyczność nie jest snopem.
Oblicz sekcje snopu stałego $k_X$ na $X$.
\wskazowka{Pokryj $X$ singletonami.}
\item \textbf{Kiełki funkcji gładkich.}
Dla $x\in\R$ opisz relację równoważności na parach $(U,f)$,
$U\ni x$, $f\in C^\infty(U)$, która wyznacza włókno snopu
funkcji gładkich. Pokaż, że wartości $f(x)$ nie wystarczają
do rozróżniania wszystkich kiełków.
\item \textbf{Przedłużenie przez zero.}
Dla $j:(0,1)\hookrightarrow\R$ uzasadnij na podstawie definicji,
że $(j_!k_{(0,1)})_0=0$, a $(j_*k_{(0,1)})_0\cong k$.
Wyjaśnij, dlaczego lokalnie stała funkcja równa $1$ na $(0,1)$
nie wyznacza globalnej sekcji $j_!k_{(0,1)}$.
\item \textbf{Mayer--Vietoris dla nośników.}
Dla otwartego pokrycia $X=U\cup V$ sprawdź dokładność na
włóknach ciągu
$0\to (j_{U\cap V})_!\mathcal F\to (j_U)_!\mathcal F\oplus(j_V)_!\mathcal F\to\mathcal F\to0$.
Dlaczego kolejność strzałek jest odwrotna od zwykłego ciągu
\v{C}echa sekcji na $X$?
\item \textbf{Lokalny system na okręgu.}
Niech $A=\begin{pmatrix}1&1\\0&1\end{pmatrix}$ działa na $k^2$.
Oblicz $H^0(S^1;L_A)$ i $H^1(S^1;L_A)$, wskazując
konkretny generator każdego jednowymiarowego wyniku.
\wskazowka{Zbadaj jądro i iloraz przez obraz $A-I$.}
\item \textbf{Nośniki zwarte i de Rham.}
Dla $M=\R$ udowodnij, że jedyna gładka funkcja lokalnie
stała o nośniku zwartym jest zerowa, a klasa
$\rho(t)\,\mathrm dt$, $\int_{t\in\R}\rho(t)\,\mathrm dt=1$,
generuje $H_{c,\mathrm{dR}}^1(\R)$.
Wskaż, jak wynik wynika z rezolwenty drobnej.
\item \textbf{Kontrola przesunięcia.}
Załóż, że $V$ jest kompleksem skoncentrowanym w stopniu $0$.
W którym stopniu znajduje się jedyna grupa kohomologii $V[3]$?
Oblicz $H^{-2}(\R^2,k_{\R^2}[2])$.
\item \textbf{Węzeł krzywej.}
Dla $X=\{zw=0\}\cap B_\varepsilon\subset\C^2$ oblicz
$H^1(X,X\setminus\{0\};k)$ i $H^2(X,X\setminus\{0\};k)$.
Porównaj z wynikiem dla gładkiego punktu krzywej zespolonej.
\wskazowka{Usuń węzeł i sprowadź problem do dwóch rozłącznych okręgów.}
\item \textbf{Właściwość i obraz bezpośredni.}
Wyjaśnij, dlaczego dla właściwej mapy $f:X\to Y$
warunek właściwości nośnika sekcji względem $Y$ jest
spełniony automatycznie. Podaj przykład zanurzenia otwartego,
dla którego $f_!$ i $f_*$ nie są izomorficzne.
\item \textbf{Poincaré jako szczególny przypadek Verdiera.}
Dla spójnej orientowalnej rozmaitości $M$ wymiaru $n$
przyjmij $\omega_M^\bullet\simeq k_M[n]$ i zapisz
wzór dualności Verdiera dla $\mathcal F=k_M$.
Wyjaśnij, skąd bierze się znak przesunięcia stopnia
w dualności $H_c^q(M;k)^\vee\cong H^{n-q}(M;k)$.
\wskazowka{Zastosuj $H^{-q}(M,k_M[n])=H^{n-q}(M;k)$.}
\item \textbf{Ciąg lokalizacji.}
Dla $X=\R$, $Z=\{0\}$ i $U=\R\setminus\{0\}$
wyprowadź dokładny ciąg
$0\to k\to k^2\to k\to0$ z lokalizacji nośników.
Wyznacz mapy przy spójnym wyborze orientacji półprostych.
\item \textbf{Wiązki $\mathcal O(m)$ na $\mathbb P^1$.}
Korzystając z dwóch standardowych kart prostej rzutowej,
oblicz $H^0(\mathbb P^1,\mathcal O(2))$ i
$H^1(\mathbb P^1,\mathcal O(-4))$.
Wskaż reprezentantów bazowych w kompleksie \v{C}echa.
\wskazowka{Dla $\mathcal O(-4)$ pozostają wykładniki
$-3,-2,-1$, których nie usuwają kobrzegi.}
\end{enumerate}

\chapter{Homologia przecięć, snopy perwersyjne i twierdzenie o rozkładzie}
\begin{enumerate}
\item \textbf{Dopuszczalność i kowymiar.}
Dla pseudorozmaitości wymiaru $4$ z izolowanym punktem
osobliwym ustal $\bar m(4)=1$. Dla których $i$
łańcuch $i$-wymiarowy może spotkać osobliwość?
Osobno sprawdź warunek dla jego brzegu.
\wskazowka{Porównaj $0$ z $i-4+\bar m(4)$ i $(i-1)-4+\bar m(4)$.}

\item \textbf{Dwie środkowe perwersyjności.}
Wykaż, że $\bar m(c)+\bar n(c)=c-2$.
Dla których kowymiarów $c$ te funkcje są równe?

\item \textbf{Stożek na sferze.}
Dla $L=S^{n-1}$ sprawdź, że wzór stożkowy
odtwarza grupy homologii $\R^n$ (dla $n\ge2$).
Czy wynik zależy od $\bar p$?
\wskazowka{Jedyna dodatnia homologia $L$ występuje w stopniu $n-1$.}

\item \textbf{Stożek na torusie.}
Oblicz $IH_i^{\bar m}(cT^2;\Q)$ i
$IH_i^{\bar n}(cT^2;\Q)$.
Wyprowadź próg dopuszczalności bez odwoływania się
do gotowego wyniku dla stożka.

\item \textbf{Zawieszenie torusa.}
Dla $X=\Sigma T^2$ wykonaj pełny rachunek długiego
ciągu Mayera--Vietorisa dla $\bar m$ i $\bar n$,
uwzględniając homomorfizmy w stopniu zerowym.
Sprawdź, które pary grup uczestniczą w dualności.
\wskazowka{W stopniu pierwszym mapa dla $\bar m$ jest,
z dokładnością do znaku, przekątną $k^2\to k^2\oplus k^2$.}

\item \textbf{Snop perwersyjny na gładkiej krzywej.}
Niech $M$ będzie gładką spójną krzywą zespoloną.
Policz stopnie niezerowych grup
$H^j(i_x^*k_M[1])$ i $H^j(i_x^!k_M[1])$.
Sprawdź oba warunki definicji perwersyjności.

\item \textbf{Normalizacja węzła.}
Dla $X=\{zw=0\}\cap B_\varepsilon$ i jej normalizacji
$\nu$ policz wszystkie niezerowe grupy kohomologii
kiełka i kopiełka $\nu_*k_{\widetilde X}[1]$ w węźle.
Dlaczego nie może istnieć podobiekt ani iloraz
perwersyjny podparty wyłącznie w węźle?
\wskazowka{Użyj dwóch dysków i lokalnej kohomologii pary
$(D^2,D^2\setminus\{0\})$.}

\item \textbf{Węzłowa krzywa rzutowa.}
Załóż, że $C$ jest krzywą otrzymaną z $\mathbb P^1$
przez utożsamienie dwóch różnych punktów.
Oblicz $H^i(C;\Q)$ i $IH^i(C;\Q)$.
Wyjaśnij, skąd bierze się różnica w stopniu $1$.

\item \textbf{Małość rozdzielczości.}
Sprawdź, że normalizacja krzywej węzłowej jest mała,
a rozdmuchanie punktu na $\C^2$ jest półmałe,
lecz nie małe. Wskaż dokładne nierówności wymiarowe.

\item \textbf{Rozdmuchanie i pochodne obrazy.}
Niech $f:\widetilde{\C^2}\to\C^2$ oznacza rozdmuchanie
początku. Oblicz wszystkie $R^qf_*\Q$ przy użyciu
włókien mapy. Dlaczego sama ta informacja nie jest
ogólnie wystarczająca do rozszczepienia $Rf_*\Q$?
\wskazowka{Oddziel obliczenia kiełków od problemu klas rozszerzeń.}

\item \textbf{Przesunięcia w formule rozkładu.}
Dla $f:\mathbb P^1\times\mathbb P^1\to\mathbb P^1$
zapisz $Rf_*\Q[2]$ jako sumę przesunięć
perwersyjnego snopu $P=\Q_{\mathbb P^1}[1]$.
Wyznacz jego niezerowe ${}^pH^j$.

\item \textbf{Porównanie zakresów twierdzeń.}
Wypisz hipotezy gwarantujące istnienie
(i) łańcuchowej homologii przecięć,
(ii) dualności Goresky'ego--MacPhersona,
(iii) twierdzenia o rozkładzie projektowym.
Podaj powód, dla którego nie wolno automatycznie
zastępować w (iii) „projektowy” słowem „ciągły”.
\end{enumerate}

\chapter{Kategorie pochodne, Hodge, cykle zanikające i monodromia}
\begin{enumerate}
\item \textbf{Stożek i znaki.}
Dla dowolnego morfizmu kompleksów $f:A^\bullet\to B^\bullet$
wykaż bezpośrednio, że $d_{\operatorname{Cone}(f)}^2=0$.
Sprawdź wzór na homomorfizm łączący ciągu stożka.
\wskazowka{Zastosuj $d(b,a)=(d_Bb+f(a),-d_Aa)$ i użyj $d_Bf=fd_A$.}

\item \textbf{Mnożenie przez liczbę.}
Oblicz kohomologię stożka $\Z[0]\xrightarrow{\cdot m}\Z[0]$
dla $m=0$, $m=1$ i $|m|\ge2$.
Sformułuj warunek na quasi-izomorfizm.

\item \textbf{Filtracja czystej struktury Hodge'a.}
Dla czystej struktury Hodge'a wagi $m$ wyprowadź
$V_\C=F^p\oplus\overline{F^{m-p+1}}$.
Odtwórz wszystkie $V^{p,q}$ z samej filtracji $F$.
\wskazowka{Wyznacz wspólne składniki dwóch filtracji.}

\item \textbf{Diament Hodge'a torusa zespolonego.}
Dla krzywej eliptycznej $E$ policz wszystkie $h^{p,q}$.
Sprawdź $b_1(E)=2$ i $\chi(E)=0$.

\item \textbf{Nieparzysta liczba Bettiego.}
Udowodnij, że gdy zwarta rozmaitość zespolona ma
$b_1$ nieparzyste, nie może być Kählerowska.
Zastosuj wynik do powierzchni Hopfa, wiedząc, że
$S^1\times S^3$ jest jej modelem topologicznym.
\wskazowka{Sparuj $H^{p,q}$ z $H^{q,p}$.}

\item \textbf{Obliczenie wagi węzła.}
Przepisz dokładny ciąg snopów normalizacji
$0\to\Q_C\to\nu_*\Q_{\widetilde C}\to\Q_{\{x\}}\to0$.
Oblicz długi ciąg kohomologii i udowodnij,
że $H^1(C;\Q)$ ma wagę $0$.
\wskazowka{Globalne funkcje stałe dają zerową różnicę wartości gałęzi.}

\item \textbf{Osobliwość węzłowa.}
Dla $f(z,w)=zw$ w $\C^2$ opisz włókno Milnora
jako pierścień i wyznacz zredukowane kohomologie.
Sprawdź monodromię na $H^1$ przez jawny transport.

\item \textbf{Forma kwadratowa i model $T^*S^n$.}
Wykaż, że afiniczna kwadryka $\sum z_i^2=1$
jest dyfeomorficzna z $T^*S^n$. Zapisz odwrotność
i skonstruuj deformacyjną retrakcję na $S^n$.
\wskazowka{Zapisz $z=x+iy$ i wyprowadź $x\perp y$ oraz $|x|^2=1+|y|^2$.}

\item \textbf{Parzystość i monodromia.}
Dla $n=1,2,3,4$ policz działanie monodromii
osobliwości $\sum_{j=0}^nz_j^2$ na
$\widetilde H^n(F_t;\Z)$.
Dlaczego sama liczba cykli zanikających nie wystarcza
do odróżnienia tych przypadków?

\item \textbf{Transwekcja Picarda--Lefschetza.}
Załóż $\langle b,a\rangle=1$ w $H_1(E;\Z)$
i oblicz macierz monodromii dla cyklu $a$.
Sprawdź, że $T$ zachowuje formę przecięcia.
\wskazowka{Wylicz bezpośrednio $T^tJT$ dla odpowiedniej macierzy $J$.}

\item \textbf{Cykle pobliskie i zanikające.}
Wyprowadź długi ciąg grup kohomologii kiełków
z trójkąta $i^*K\to R\psi_fK\to R\phi_fK\to i^*K[1]$.
W przypadku $f(z,w)=zw$, $K=\Q_X$ wskaż,
które stopnie mogą być niezerowe.

\item \textbf{Obraz prosty a jego kohomologie.}
Wyjaśnij, dlaczego dla $f:Y\to X$ znajomość
wszystkich $R^qf_*\Q_Y$ nie musi wyznaczać
$Rf_*\Q_Y$. Wypisz, jakiego dodatkowego wyniku
potrzeba do rozkładu dla morfizmu projektowego
z gładkiej odmiany $Y$.
\wskazowka{Zbadaj trójkąty truncacji i możliwe morfizmy łączące.}

\item \textbf{Porównanie teorii.}
Dla nieredukowalnej krzywej wymiernej z jednym węzłem
zestaw $H^1(C;\Q)$, $IH^1(C;\Q)$, wagę $H^1(C)$
i lokalną grupę cykli zanikających przy wygładzaniu węzła.
Napisz, które obiekty są globalne, a które lokalne.
\end{enumerate}

\chapter{Wariacje struktur Hodge'a, okresy i degeneracje}
\begin{enumerate}
\item \textbf{Lokalny system i płaska koneksja.}
Niech $S=\C^*$, a $\rho:\pi_1(S)\to\mathrm{GL}_2(\Z)$
wyznacza macierz $T=\begin{psmallmatrix}1&m\\0&1\end{psmallmatrix}$.
Zbuduj odpowiadający lokalny system, jego płaską wiązkę
oraz opis przekrojów poziomych na uniwersalnym nakryciu.
\wskazowka{Użyj ilorazu $(\C\times\C^2)/(z,v)\sim(z+1,Tv)$ i sprawdź, że na nakryciu koneksją jest zwykła różniczka.}

\item \textbf{Transwersalność nie znaczy poziomości.}
Dla filtracji $F^1\subset F^0=\mathcal H$ struktury Hodge'a wagi jeden
wyjaśnij, dlaczego warunek Griffithsa nie wymusza
$\nabla(F^1)\subseteq F^1\otimes\Omega_S^1$.
\wskazowka{Dla $p=1$ prawa strona warunku Griffithsa to całe $F^0=\mathcal H$. Rozważ nieizotrwialną rodzinę krzywych eliptycznych.}

\item \textbf{Całka beta.}
Dla $n\ge0$ wyprowadź
$\int_0^1 s^{n-1/2}(1-s)^{-1/2}\,\dd s
=\pi\binom{2n}{n}/4^n$.
\wskazowka{Podstaw $s=\sin^2\theta$ i zastosuj rekurencję całek potęg funkcji cosinus.}

\item \textbf{Równanie hipergeometryczne.}
Dla $a_n=\binom{2n}{n}^2/16^n$ sprawdź
$(n+1)^2a_{n+1}=(n+1/2)^2a_n$.
Wyprowadź równanie Picarda--Fuchsa oraz znajdź pierwsze
cztery wyrazy rozwinięcia $u(\lambda)$.
\wskazowka{Porównaj współczynniki przy kolejnych potęgach $\lambda$ po podstawieniu szeregu do równania.}

\item \textbf{Osobliwości równania Picarda--Fuchsa.}
Znajdź wykładniki indycjalne przy $\lambda=0$ i
$\lambda=1$. Dlaczego drugie lokalne rozwiązanie
może zawierać logarytm?
\wskazowka{Podstaw $u(\lambda)=\lambda^r$ do najniższych rzędów równania; obydwa pierwiastki równania indycjalnego są równe zeru.}

\item \textbf{Dlaczego monodromia Legendre'a ma wykładnik dwa?}
Sprawdź lokalną holomorficzną zmianę współrzędnych
$Y=y/\sqrt{1-x}$, $X=x-\lambda/2$ dla
$y^2=x(x-1)(x-\lambda)$. Wyprowadź $UV=\lambda^2/4$
oraz porównaj działanie monodromii z modelem $UV=t$.
\wskazowka{Wybierz gałąź pierwiastka holomorficznego blisko $x=0$, a następnie $U=Y+\mathrm iX$, $V=Y-\mathrm iX$.}

\item \textbf{Logarytm nilpotentnego operatora.}
Niech $T=\begin{psmallmatrix}1&m\\0&1\end{psmallmatrix}$ z $m\ne0$.
Oblicz $N=\log T$, $\ker N$ i $\operatorname{im}N$.
Wyznacz $W_\bullet$ i sprawdź, że $N$ indukuje izomorfizm
$\operatorname{Gr}_2^W\to\operatorname{Gr}_0^W$.
\wskazowka{Mamy $(T-I)^2=0$, więc szereg logarytmu kończy się po pierwszym wyrazie.}

\item \textbf{Włókno centralne a struktura graniczna.}
Niech $C_0$ będzie nieredukowalną krzywą wymierną z jednym węzłem,
a $E_t$ jej wygładzaniem o nietrywialnej unipotentnej monodromii.
Porównaj wymiar i rozkład wag na $H^1(C_0;\Q)$ oraz na
$H^1(E_t;\Q)$ wyposażonym w graniczną mieszaną strukturę Hodge'a.
\wskazowka{Dla $C_0$ użyj normalizacji; dla struktury granicznej oblicz $\operatorname{im}N=\ker N$ w wymiarze dwa.}

\item \textbf{Kontrola orientacji i macierzy monodromii.}
Jeżeli $u\mapsto u$, $v\mapsto v+2u$, wyznacz macierz
na kolumnie $(u,v)^t$ oraz działanie na ilorazie $\tau=v/u$.
Wyjaśnij, dlaczego macierz operatora na bazie funkcji $(u,v)$
jest transpozycją macierzy działającej na kolumnę wartości
$(u(\lambda),v(\lambda))^t$.
\wskazowka{Uważaj na różnicę między działaniem na funkcjach okresowych, na wektorach ich współrzędnych i na cyklach.}
\end{enumerate}


\chapter{Sześć operacji snopowych, dualność i cykle zanikające}
\begin{enumerate}
\item \textbf{Stożek morfizmu.}
Dla morfizmu kompleksów $u:A\to B$ sprawdź $d^2=0$ w $\operatorname{Cone}(u)$,
a następnie wyprowadź długi ciąg jego kohomologii.
\wskazowka{Użyj krótkiego ciągu kompleksów $0\to B\to\operatorname{Cone}(u)\to A[1]\to0$ i zapisz morfizm łączący.}

\item \textbf{Obraz bezpośredni a przedłużenie przez zero.}
Dla $j:(0,1)\hookrightarrow\R$ wyznacz kiełki $j_!k_{(0,1)}$
oraz $j_*k_{(0,1)}$ w punktach $0,\frac12,1$.
\wskazowka{W każdym punkcie zbadaj małe przedziały otwarte i warunek na nośnik.}

\item \textbf{Lokalizacja.}
Niech $i:Z\hookrightarrow X$ będzie włożeniem zamkniętym,
$j:X\setminus Z\hookrightarrow X$ otwartym. Dowiedź dokładności
$0\to j_!j^*\mathcal F\to\mathcal F\to i_*i^*\mathcal F\to0$.
\wskazowka{Sprawdź osobno kiełki w $Z$ i w $X\setminus Z$.}

\item \textbf{Kohomologia z nośnikiem w punkcie.}
Dla $i:\{0\}\hookrightarrow\R^n$ oblicz $i^!k_{\R^n}$.
\wskazowka{Użyj $H^q(\R^n,\R^n\setminus\{0\};k)$ i deformacyjnej retrakcji do $S^{n-1}$. Wynikiem jest $k[-n]$.}

\item \textbf{Dualność na prostej.}
Wykaż $H_c^0(\R;k)=0$, $H_c^1(\R;k)\cong k$ i
sprawdź $\omega_{\R}\simeq k_{\R}[1]$.
\wskazowka{Porównaj z kohomologią pary $([0,1],\{0,1\})$.}

\item \textbf{Projekcja skończonej rodziny punktów.}
Niech $f:\{1,\dots,r\}\times Y\to Y$ będzie projekcją.
Oblicz $Rf_*\mathcal F$ i $Rf_!\mathcal F$ dla lokalnych systemów
na poszczególnych składowych. Sprawdź zmianę podstawy.
\wskazowka{Włókno jest skończone i dyskretne; wyższe grupy kohomologii znikają.}

\item \textbf{Właściwe odwzorowanie kwadratowe.}
Dla $f:\R\to\R$, $f(x)=x^2$, opisz $R^qf_*k_{\R}$ i
wyjaśnij, dlaczego $R^0f_*k_{\R}$ nie jest lokalnym systemem.
\wskazowka{Oblicz kohomologię pustego, jedno- i dwupunktowego włókna.}

\item \textbf{Wzór projekcji dla lokalnego systemu.}
Wyprowadź $Rf_!(A\otimes f^*L)\cong Rf_!A\otimes L$ dla
lokalnego systemu $L$ o skończonym rzędzie.
\wskazowka{Zredukuj lokalnie do $L=k^r$, a następnie sprawdź sklejenie naturalnych izomorfizmów.}

\item \textbf{Cykle zanikające odwzorowania $z^m$.}
Dla $f:\C\to\C$, $f(z)=z^m$, $m\ge2$, wyznacz
$(R^0\phi_fk_{\C})_0$ i monodromię jako permutację $m$ korzeni.
\wskazowka{Włókno Milnora ma $m$ punktów; iloraz $k^m/\Delta k$ ma wymiar $m-1$.}

\item \textbf{Osobliwość węzłowa.}
Dla $f(z,w)=zw$ oblicz $R\psi_f$ i $R\phi_f$ w punkcie $0$.
Wyjaśnij różnicę między monodromią lokalną a działaniem Picarda--Lefschetza
na globalnej homologii gładkiej krzywej projektowej.
\wskazowka{Lokalne włókno Milnora jest homotopijnie równoważne $S^1$; jego monodromia na $H^1$ jest identycznością.}

\item \textbf{Wymiana $*$ i $!$ przez dualność.}
Przyjmij formalizm sześciu operacji i wyprowadź
$\mathbb D_Y Rf_!\simeq Rf_*\mathbb D_X$.
Zaznacz, gdzie używasz adjunkcji i wzoru projekcji.
\wskazowka{Przetestuj obydwie strony morfizmami z dowolnego obiektu $B$, a następnie użyj lematu Yonedy.}

\item \textbf{Kierunek implikacji przy zmianie podstawy.}
Wyjaśnij, dlaczego przy niewłaściwym $f$ nie należy automatycznie
utożsamiać $(R^qf_*A)_y$ z $H^q(f^{-1}(y);A)$.
\wskazowka{Właściwa zmiana podstawy jest twierdzeniem z dodatkowymi hipotezami; rozważ zachowanie rodzin niezwartych włókien.}
\end{enumerate}

\chapter{Topologia odmian projektowych: Lefschetz, formy przecięcia i rozdmuchania}
\begin{enumerate}
\item \textbf{Słabe twierdzenie Lefschetza w stopniach granicznych.}
Niech $X\subset\mathbb P^N_\C$ będzie gładką odmianą projektową
wymiaru $n\ge2$, a $Y$ gładkim przekrojem hiperpłaszczyznowym.
Wyprowadź z zanikania $H^q(X,Y;\Z)$ dla $q<n$ zakres
izomorfizmów i monomorfizmu $H^q(X;\Z)\to H^q(Y;\Z)$.
\wskazowka{Zapisz trzy kolejne wyrazy długiego ciągu kohomologii pary.}

\item \textbf{Zanikające cykle i grupa podstawowa.}
Dla gładkiej powierzchni projektowej $X$ i gładkiej krzywej
hiperpłaszczyznowej $Y$ wykaż surjektywność $\pi_1(Y)\to\pi_1(X)$.
Czy można bez dodatkowych hipotez zastąpić ją izomorfizmem?
\wskazowka{Użyj ciągu dokładnego pary $(X,Y)$ i porównaj $Y$ będącą krzywą o dodatnim rodzaju z $X=\mathbb P^2$.}

\item \textbf{Lemat o łańcuchach Lefschetza.}
Uzasadnij indukcyjnie $\Lambda L^rv=r(n-k-r+1)L^{r-1}v$
dla $v\in\ker\Lambda\cap V^k$ oraz wyprowadź formułę na normy.
\wskazowka{Z komutatora $[\Lambda,L]=H$ i równości $HL=L(H-2)$ wyznacz $[\Lambda,L^r]$.}

\item \textbf{Liczby Bettiego i Hodge'a.}
Wykaż, że $b_{2j+1}(X)$ jest parzyste dla każdej gładkiej
spójnej odmiany projektowej $X$ nad $\C$.
\wskazowka{W rozkładzie Hodge'a dla nieparzystego stopnia nie ma składnika diagonalnego $H^{p,p}$.}

\item \textbf{Dodatniość i ujemność na powierzchni.}
Niech $S$ będzie gładką projektową powierzchnią i $\kappa$ klasą Kählerowską.
Wykaż, że forma przecięcia jest ujemnie określona
na $\{x\in H^{1,1}(S;\R):\int_S x\wedge\kappa=0\}$.
\wskazowka{W stopniu drugim prymitywność znaczy $Lx=0$; użyj relacji Hodge'a--Riemanna.}

\item \textbf{Formy przecięcia trzech powierzchni.}
Oblicz sygnatury $\mathbb P^2_\C$, $\mathbb P^1_\C\times\mathbb P^1_\C$
oraz $\operatorname{Bl}_p\mathbb P^2_\C$.
\wskazowka{Na iloczynie weź klasy $a,b$ czynników; zachodzi $a^2=b^2=0$, $\int a\smile b=1$.}

\item \textbf{Rozdmuchanie w punkcie i klasy Cherna.}
Dla $\pi:\widetilde S\to S$ wyprowadź formuły
$K_{\widetilde S}^2=K_S^2-1$ oraz
$\chi(\widetilde S)=\chi(S)+1$.
\wskazowka{W pierwszej formule użyj $K_{\widetilde S}=\pi^*K_S+E$, $E^2=-1$; w drugiej zastąp punkt wyjątkowym $\mathbb P^1$.}

\item \textbf{Wiązka normalna.}
W dwóch kartach lokalnego rozdmuchania $(x,y)$ punktu
w $\C^2$ zapisz przejścia $u=y/x$, $v=x/y$ i
udowodnij $N_{E/\widetilde S}\cong\mathcal O_{\mathbb P^1}(-1)$.
\wskazowka{Śledź zmianę współrzędnej \emph{normalnej}, a nie stycznej do $E$.}

\item \textbf{Nie-Kählerowskość klasy ściągniętej.}
Wyjaśnij, dlaczego $\pi^*h$ nie może być klasą
Kählerowską na $\operatorname{Bl}_p\mathbb P^2$, choć
$\int_{\widetilde S}(\pi^*h)^2>0$.
\wskazowka{Policz całkę $\int_E\pi^*h$ po krzywej wyjątkowej.}

\item \textbf{Rozdmuchanie w krzywej.}
Niech $X$ będzie trójwymiarową odmianą projektową,
a $Z\subset X$ gładką krzywą rodzaju $g$.
Wyznacz zmiany $b_2$, $b_3$, $b_4$ po rozdmuchaniu w $Z$.
\wskazowka{Dla kodymensji $r=2$ nowa część to $H^{k-2}(Z;\Q)$.}

\item \textbf{Przykład przecięcia torusowego.}
Oblicz macierz przecięcia na $H^2(\mathbb P^1\times\mathbb P^1;\Z)$
i wyjaśnij jej zgodność z twierdzeniem o indeksie Hodge'a.
\wskazowka{W bazie klas hiperpłaszczyzn pochodzących z dwóch czynników otrzymasz macierz z zerami na przekątnej.}

\item \textbf{Rozróżnienie twierdzeń Lefschetza.}
Sformułuj osobno słabe, silne i względne silne twierdzenie Lefschetza.
Podaj przestrzenie, obiekty, mapy i zakres stopni w każdej tezie.
\wskazowka{Porównaj restrykcję na przekrój, mnożenie przez klasę Kählerowską i działanie w perwersyjnej kohomologii obrazu właściwego morfizmu.}
\end{enumerate}

\chapter{Granice odwrotne, wyższe przeszkody i dalsza teoria homotopii}
\begin{enumerate}
\item \textbf{Rozwiązanie równań różnicowych.}
Dla układu $A_j=\Z/2^j\Z$ z redukcjami pokaż, że
$\varprojlim^1A_j=0$. Wyznacz $\varprojlim A_j$.
\wskazowka{Mapy przejścia są surjektywne; granica to pierścień liczb $2$-adycznych.}
\item \textbf{Obliczenie $\varprojlim^1$.}
Dla układu $\Z\xleftarrow{\cdot3}\Z\xleftarrow{\cdot3}\cdots$
wykaż, że $\varprojlim A_j=0$ i $\varprojlim^1A_j\cong\Z_3/\Z$.
\wskazowka{Zbuduj mapę $(a_j)\mapsto\sum 3^ja_j$ i określ jej jądro po przejściu do ilorazu.}
\item \textbf{Produkty i homologia.}
Dla układu kompleksów, w którym mapy są surjektywne w każdym stopniu,
wyprowadź ciąg Milnora z ciągu dokładnego kompleksów
$0\to\varprojlim C^{(j)}\to\prod C^{(j)}\xrightarrow{1-T}\prod C^{(j)}\to0$.
\wskazowka{Sprawdź najpierw, że $H_n(\prod_j C^{(j)})\cong\prod_jH_n(C^{(j)})$.}
\item \textbf{Uwaga na granicę przestrzeni.}
Wyjaśnij, dlaczego z homeomorfizmu $X=\varprojlim X_j$ nie wynika
izomorfizm $C_*(X)\cong\varprojlim C_*(X_j)$.
\wskazowka{Pojedynczy sympleks singularny w jednym z $X_j$ może nie mieć zgodnych podniesień.}
\item \textbf{Wyższa przeszkoda na $S^4$.}
Wykaż, że każda zorientowana wiązka rzeczywista rzędu $3$ nad $S^4$
ma zerową klasę Eulera, lecz nie każda ma przekrój nigdzie niezerowy.
\wskazowka{Sklejanie klasyfikuje takie wiązki przez $\pi_3(SO(3))$, a przekrój sprowadza grupę strukturalną do $SO(2)$.}
\item \textbf{Zmiana wyboru częściowego przekroju.}
Załóż, że pierwsza przeszkoda jest reprezentowana przez kołańcuch
$o_{q+1}$, a zmianę przekroju na $q$-komórkach koduje $\beta$.
Wyjaśnij, czemu nowa przeszkoda jest $o_{q+1}+\delta\beta$.
\wskazowka{Śledź, jak mapa na brzegu $(q+1)$-dysku zmienia się pod wpływem modyfikacji na jego $q$-ścianach.}
\item \textbf{Relacje Adema w małych stopniach.}
Wyprowadź z wzoru Adema równości
$\operatorname{Sq}^1\operatorname{Sq}^2=\operatorname{Sq}^3$ i
$\operatorname{Sq}^2\operatorname{Sq}^2=\operatorname{Sq}^3\operatorname{Sq}^1$.
\wskazowka{Wystarczy podstawić $(a,b)=(1,2)$ i $(2,2)$ oraz policzyć dwumiany modulo dwa.}
\item \textbf{Kontrola na przestrzeni rzutowej.}
Oblicz $\operatorname{Sq}^i(t^m)$ na $H^*(\mathbb RP^\infty;\F_2)$,
a następnie dla $(a,b)=(2,2)$ sprawdź relację Adema na
wszystkich jednomianach $t^m$.
\wskazowka{Korzystaj z $\operatorname{Sq}(t)=t+t^2$ i tożsamości dwumianowych modulo dwa. Sama zgodność nie jest dowodem relacji naturalnych.}
\item \textbf{Blakers--Massey i Freudenthal.}
Załóż twierdzenie Blakersa--Masseya i udowodnij, że dla
$(r-1)$-spójnego CW-kompleksu $Y$, $r\ge2$, mapa
$\pi_k(Y)\to\pi_{k+1}(\Sigma Y)$ jest izomorfizmem dla
$k\le2r-2$ i surjekcją dla $k=2r-1$.
\wskazowka{Rozłóż zawieszenie na dwa stożki i porównaj dwie grupy homotopii względnej.}
\item \textbf{Reszta przy sklejaniu trajektorii.}
Wyjaśnij, dlaczego obcięcie trajektorii o wykładniczym zbiegu
do punktu krytycznego daje błąd $O(e^{-cR})$, a małość błędu
nie zapewnia jeszcze dokładnej trajektorii.
\wskazowka{Oszacowanie residuum jest elementarne; do równania korekcyjnego potrzebny jest jednostajnie ograniczony prawy odwrotny operatora linearyzacji.}
\item \textbf{Status dowodu.}
Sformułuj trzy osobne twierdzenia potrzebne do uzyskania
$\partial^2=0$ w kompleksie Morse'a: zwartość do łamania,
lokalne sklejanie i zgodność orientacji. Wyjaśnij zależności między nimi.
\wskazowka{Brzeg zwartej zorientowanej jednowymiarowej przestrzeni moduli jest podpisaną sumą złamanych trajektorii.}
\end{enumerate}

\end{document}
